\documentclass[11pt]{report}
\usepackage{fontspec}
\setmainfont{TeX Gyre Pagella}
\usepackage[margin=1.15in]{geometry}
\usepackage{amsmath,amssymb,amsthm}
\usepackage{enumitem}
\usepackage{hyperref}
\hypersetup{colorlinks=true, linkcolor=black, urlcolor=black, citecolor=black}

\newtheorem{proposition}{Proposition}[chapter]
\newtheorem{theorem}{Theorem}[chapter]
\newtheorem{lemma}{Lemma}[chapter]
\newtheorem{definition}{Definition}[chapter]
\newtheorem{remark}{Remark}[chapter]
\newtheorem{corollary}{Corollary}[chapter]
\newtheorem{example}{Example}[chapter]
\newtheorem{conjecture}{Conjecture}[chapter]
\newtheorem{openproblem}{Open Problem}[chapter]

\title{\textbf{Generative Continuation Geometry}}
\author{Flyxion}
\date{July 2026}

\begin{document}

\maketitle

\tableofcontents

\chapter{Continuation Attractors}

A fern frond is not drawn leaf by leaf. Look closely at one pinna and
it is, near enough, a smaller fern; look closer still at one pinnule
of that pinna and it is, near enough, a smaller fern again. Nobody
specified the shape of the whole plant directly. A small number of
growth rules, applied to their own output over and over, produced
something with more detail than any of the rules themselves contain.
The shape is not stored anywhere. It is generated.

\emph{Continuation Geometry} spent its length measuring futures:
how much room a state has, how far apart two states' futures sit, when
a distinction between futures collapses to nothing. All of that took
the space of futures as given, something to be measured once it
already existed. This book asks a different question. Can a space of
futures itself be grown, the way the fern is grown, from a small
number of simple moves applied to their own output? If continuation
spaces can be generated rather than merely measured, the natural unit
of study stops being the space and becomes the small set of
operations that built it.

This opening chapter answers that question in the one case where the
answer is already, in a precise sense, classical. It is not a new
theorem. It is an old and beautiful one, due to Hutchinson, applied
here to the specific metric geometry the first book already built. The
payoff is that self-similar continuation structures, the kind
suggested by a duck, a crown, and a leaf all arising from the same
underlying rule, turn out to rest on exactly the Hausdorff-distance
machinery already in hand, with nothing new invented and nothing
smuggled in from outside.

\section*{Preliminaries: expansion is not generation}

Two facts from \emph{Continuation Geometry} are worth restating here,
briefly and as what they actually are: lemmas, not theorems, and
useful precisely because they show what generation is \emph{not}.

\begin{lemma}[Reachable expansion]
For $S_t \subseteq X$ and $\mathrm{Adj}(S_t) := \bigcup_{x \in S_t} \Phi(x)$
under a one-step reachability map $\Phi$, the sequence
$S_{t+1} := S_t \cup \mathrm{Adj}(S_t)$ satisfies $S_t \subseteq S_{t+1}$
for every $t$.
\end{lemma}
\begin{proof}
Immediate: $S_t \subseteq S_t \cup \mathrm{Adj}(S_t) = S_{t+1}$ by
definition of union.
\end{proof}

\begin{remark}
This, together with the reachable-set recursion proved by induction
in \emph{Continuation Geometry} (Kauffman's adjacent possible as an
instance of finite-horizon reachability), is the full content of what
that book called generation, and it is worth being honest that it
earns the name only loosely. The sequence $S_0, S_1, S_2, \dots$ grows
monotonically and without bound unless the underlying space is
finite; nothing about it converges, and nothing about it is
self-similar. It is accumulation, not generation in the sense the
fern illustrates. What follows is different in kind: rather than a
sequence of ever-larger unions, it is a fixed operator applied
repeatedly to itself, and the interesting behavior is not growth but
convergence to a shape that reproduces its own generating rule.
\end{remark}

\section*{The space of shapes}

Fix $X = \mathbb{R}^n$ (or any complete metric space) and let
$\mathcal{K}(X)$ denote the set of nonempty compact subsets of $X$,
equipped with the Hausdorff distance $d_H$ of Chapter 6 of
\emph{Continuation Geometry}.

\begin{theorem}[Completeness of the space of shapes]
\label{thm:complete}
If $(X,d)$ is complete, then $(\mathcal{K}(X), d_H)$ is complete.
\end{theorem}
\begin{proof}[Proof sketch]
Given a Cauchy sequence $\{A_n\}$ in $(\mathcal{K}(X), d_H)$, define
$A := \bigcap_{N \geq 1} \overline{\bigcup_{n \geq N} A_n}$. One shows,
using completeness of $X$ and the Cauchy property of $\{A_n\}$ in
$d_H$, that $A$ is nonempty and compact, and that $d_H(A_n, A) \to 0$.
This is a standard but genuinely technical result in its full form,
requiring care with total boundedness and a diagonal argument to
extract convergent selections from each $A_n$; the complete argument
is given in \cite{falconer2003} and \cite{barnsley1988} and is not
reproduced in full here.
\end{proof}

\begin{remark}
This theorem is doing real work that is easy to take for granted.
Without it, nothing built on top of $(\mathcal{K}(X), d_H)$ would be
guaranteed to have a limit at all; a sequence of shapes could get
closer and closer together, in the sense of Hausdorff distance,
without converging to any actual shape. That this cannot happen, that
the space of compact reachable sets is itself complete whenever the
underlying state space is, is the entire reason the construction
below is guaranteed to produce something rather than merely getting
closer to nothing in particular.
\end{remark}

\section*{Contractions and the Hutchinson operator}

\begin{definition}[Contraction]
$T : X \to X$ is a contraction with constant $c \in [0,1)$ if
$d(Tx, Ty) \leq c\, d(x,y)$ for all $x, y \in X$.
\end{definition}

\begin{lemma}[Contractions act contractively on shapes]
\label{lem:setcontract}
If $T$ is a contraction on $X$ with constant $c$, then for
$A, B \in \mathcal{K}(X)$, $d_H(T(A), T(B)) \leq c\, d_H(A,B)$, where
$T(A) := \{Ta : a \in A\}$.
\end{lemma}
\begin{proof}
Fix $a \in A$. Since $B$ is compact, there is $b \in B$ with
$d(a,b) = \operatorname{dist}(a,B) \leq d_H(A,B)$. Then
$d(Ta,Tb) \leq c\, d(a,b) \leq c\, d_H(A,B)$, and since
$Tb \in T(B)$, $\operatorname{dist}(Ta, T(B)) \leq c\, d_H(A,B)$.
Taking the supremum over $a \in A$ gives
$\sup_{a' \in T(A)} \operatorname{dist}(a', T(B)) \leq c\, d_H(A,B)$.
The symmetric argument, exchanging the roles of $A$ and $B$, gives the
same bound the other way, so $d_H(T(A),T(B)) \leq c\, d_H(A,B)$.
\end{proof}

\begin{lemma}[Unions are 1-Lipschitz]
\label{lem:unionlip}
For finite collections $\{A_i\}_{i=1}^n$, $\{B_i\}_{i=1}^n$ of nonempty
compact sets, $d_H\left(\bigcup_i A_i, \bigcup_i B_i\right) \leq
\max_i d_H(A_i, B_i)$.
\end{lemma}
\begin{proof}
Let $m := \max_i d_H(A_i,B_i)$. For $x \in \bigcup_i A_i$, $x \in A_j$
for some $j$, so $\operatorname{dist}(x, B_j) \leq d_H(A_j,B_j) \leq m$,
and since $B_j \subseteq \bigcup_i B_i$,
$\operatorname{dist}(x, \bigcup_i B_i) \leq \operatorname{dist}(x,B_j) \leq m$.
Taking the supremum over $x \in \bigcup_i A_i$ gives
$\sup_{x} \operatorname{dist}(x, \bigcup_i B_i) \leq m$. The symmetric
argument, exchanging the roles of the two collections, gives the same
bound the other way, so
$d_H\left(\bigcup_i A_i, \bigcup_i B_i\right) \leq m$.
\end{proof}

\begin{definition}[Hutchinson operator]
Given contractions $T_1, \dots, T_n$ on $X$ with constants
$c_1, \dots, c_n$, define $\mathcal{H} : \mathcal{K}(X) \to \mathcal{K}(X)$
by $\mathcal{H}(A) := \bigcup_{i=1}^n T_i(A)$.
\end{definition}

\begin{proposition}[$\mathcal{H}$ is a contraction on shapes]
\label{prop:hcontract}
$d_H(\mathcal{H}(A), \mathcal{H}(B)) \leq c\, d_H(A,B)$ for
$c := \max_i c_i < 1$.
\end{proposition}
\begin{proof}
By Lemma~\ref{lem:unionlip},
$d_H(\mathcal{H}(A), \mathcal{H}(B)) = d_H\left(\bigcup_i T_i(A), \bigcup_i T_i(B)\right)
\leq \max_i d_H(T_i(A), T_i(B))$. By Lemma~\ref{lem:setcontract},
$d_H(T_i(A), T_i(B)) \leq c_i\, d_H(A,B)$ for each $i$. Hence
$d_H(\mathcal{H}(A), \mathcal{H}(B)) \leq \max_i c_i\, d_H(A,B) = c\, d_H(A,B)$.
\end{proof}

\section*{The Continuation Attractor Theorem}

\begin{theorem}[Continuation Attractor Theorem, after Hutchinson]
\label{thm:attractor}
Let $T_1, \dots, T_n$ be contractions on a complete metric space $X$
with constants $c_1, \dots, c_n$, and let $c := \max_i c_i < 1$. Then
there exists a unique $A^\ast \in \mathcal{K}(X)$ with
$\mathcal{H}(A^\ast) = A^\ast$, and for any $A_0 \in \mathcal{K}(X)$,
\[
d_H\big(\mathcal{H}^k(A_0), A^\ast\big) \leq \frac{c^k}{1-c}\, d_H(A_0, \mathcal{H}(A_0)) \xrightarrow{k \to \infty} 0.
\]
\end{theorem}
\begin{proof}
By Theorem~\ref{thm:complete}, $(\mathcal{K}(X), d_H)$ is complete. By
Proposition~\ref{prop:hcontract}, $\mathcal{H}$ is a contraction on
this complete metric space. The Banach Fixed Point Theorem
\cite{banach1922} then
supplies a unique fixed point $A^\ast$ and the stated geometric
convergence rate, applied directly to $\mathcal{H}$ on
$(\mathcal{K}(X), d_H)$.
\end{proof}

\begin{remark}
No part of this proof required inventing anything: completeness of
$\mathcal{K}(X)$ is Theorem~\ref{thm:complete}, the contraction
property of $\mathcal{H}$ is Proposition~\ref{prop:hcontract} built
from two short lemmas, and the conclusion is the Banach Fixed Point
Theorem, unmodified. This is the sense in which this chapter is
mathematically conservative even though its subject, self-similar
generated shape, sounds exotic: every piece was already sitting in
the metric structure Chapter 6 of \emph{Continuation Geometry}
built for an entirely different purpose.
\end{remark}

\section*{A worked attractor}

\begin{example}[Two refusals]
Let $X = [0,1]$, and let $T_1(x) := x/3$, $T_2(x) := x/3 + 2/3$, both
contractions with constant $c = 1/3$. The attractor $A^\ast$ satisfying
$A^\ast = T_1(A^\ast) \cup T_2(A^\ast)$ is the standard middle-thirds
Cantor set.
\end{example}

\begin{remark}
Read in the vocabulary of continuation, $T_1$ and $T_2$ are the two
survivors of a repeated refusal: at every stage, whatever occupies the
middle third of the current set of admissible outcomes is discarded,
and the remaining two-thirds are each rescaled to stand in for the
whole. What survives in the limit is neither of the two pieces nor
their simple union at any finite stage, but the unique shape that
reproduces itself exactly under the same rule that built it,
uncountable, nowhere dense, and entirely determined by two contraction
maps. Whether operators named Collapse or Refuse elsewhere in this
corpus are literally instances of $T_i$ in this technical sense, maps
satisfying the contraction inequality on some genuine underlying
metric, is a claim about those specific constructions and is not
asserted here; what is established is only that \emph{if} a family of
continuation-narrowing operators can be shown to be contractions on
$(\mathcal{K}(X), d_H)$, this theorem applies to them without
modification.
\end{remark}

\section*{What this chapter does not yet claim}

This chapter establishes that self-similar continuation structures are
mathematically available, not that any particular structure this
corpus has described informally is an instance of one. Two gaps in
particular are worth naming rather than leaving implicit. First,
nothing here says which operators arising from an actual continuation
problem, as opposed to a hand-picked pair of affine maps on $[0,1]$,
are genuinely contractions on a genuine metric; that has to be checked
case by case, the way it would for any application of a fixed-point
theorem. Second, and more speculatively, the question of whether a
small, fixed family of primitive operators could generate \emph{every}
admissible continuation geometry to arbitrary Hausdorff-distance
accuracy, a continuation-geometric analogue of universal approximation,
is not addressed by this theorem and is not claimed here. It is a
genuine open question, and answering it is a natural place for this
book to go next, once the machinery in the chapters between here and
there, quotient maps and operator composition among them, has been
built on equally solid ground.

\chapter{Continuation Quotients}

Two rivers can rise on opposite sides of a mountain range, cross
entirely different terrain, pass through valleys that share nothing
in common, and still arrive at the same confluence and become, from
that point on, one river. Everything about their separate histories,
which rocks each one wore smooth, which tributaries fed it, how long
each took to descend, becomes irrelevant downstream of the merge. A
boat launched below the confluence cannot tell which branch its water
came from. Not because the information was hidden. Because it no
longer makes a difference to anything the boat will experience from
here on.

The previous chapter showed that a continuation space can be grown
from a small number of generating operators. The natural next
question is what happens when two different generating histories,
different sequences of operators, different starting points, produce
the same continuation space at the end. When does that happen, and
what, if anything, does it mean? The temptation is to reach
immediately for language like information loss or memory collapse,
language already used loosely elsewhere in this corpus. This chapter
resists that temptation on purpose. It defines the relevant object
first, proves what can be proved about it without reference to any
interpretation, and only afterward asks what, if anything, has been
lost.

\section*{The continuation quotient}

\begin{definition}[Continuation quotient]
Let $H$ be a set of histories and let $\Gamma : H \to \mathcal{K}(X)$
assign to each history its continuation space (or, following Chapter 6
of \emph{Continuation Geometry}, its finite-horizon reachable set
$R_T$ for a fixed horizon $T$). Define $h_1 \sim_\Gamma h_2$ iff
$\Gamma(h_1) = \Gamma(h_2)$. This is an equivalence relation
(reflexivity, symmetry, and transitivity all following immediately
from equality of sets), and the continuation quotient is
$H/{\sim_\Gamma}$, with quotient map $q : H \to H/{\sim_\Gamma}$
sending $h$ to its equivalence class $[h]$.
\end{definition}

Nothing has been said yet about distinctions, memory, or loss. This is
deliberate: the object exists, and has structure, independently of any
of those interpretations.

\begin{proposition}[Well-defined continuation projection]
\label{prop:universal}
There exists a unique map $\bar\Gamma : H/{\sim_\Gamma} \to \mathcal{K}(X)$
such that $\Gamma = \bar\Gamma \circ q$.
\end{proposition}
\begin{proof}
Define $\bar\Gamma([h]) := \Gamma(h)$. This is well-defined: if
$[h_1] = [h_2]$, then $h_1 \sim_\Gamma h_2$, so
$\Gamma(h_1) = \Gamma(h_2)$ by definition of $\sim_\Gamma$, so the
value assigned to $[h_1] = [h_2]$ does not depend on which
representative was chosen. By construction,
$\bar\Gamma(q(h)) = \bar\Gamma([h]) = \Gamma(h)$ for every $h \in H$,
so $\Gamma = \bar\Gamma \circ q$. For uniqueness, suppose
$\bar\Gamma'$ also satisfies $\Gamma = \bar\Gamma' \circ q$. Since $q$
is surjective, every element of $H/{\sim_\Gamma}$ is $q(h)$ for some
$h$, and $\bar\Gamma'(q(h)) = \Gamma(h) = \bar\Gamma(q(h))$ for every
such $h$, so $\bar\Gamma' = \bar\Gamma$.
\end{proof}

\begin{remark}
This is routine, and it is worth saying plainly that it is routine:
it is the universal property of a quotient set, applied to nothing
more exotic than the map $\Gamma$ itself. What makes it worth stating
is not its difficulty but what it licenses: from here on,
$\bar\Gamma$ is a genuine function on the quotient, not merely a
suggestive notation, and anything proved about $\bar\Gamma$ is a fact
about the quotient set itself, not about any particular history used
to reach a point in it.
\end{remark}

\section*{A structural lemma}

The worked example below relies on the following fact about
time-homogeneous, deterministic-image dynamics, worth isolating here.

\begin{lemma}[One-step agreement propagates]
\label{lem:propagate}
Suppose $R_1(x) = R_1(y)$ for states $x, y$ under a time-homogeneous
reachability map. Then $R_T(x) = R_T(y)$ for every $T \geq 1$.
\end{lemma}
\begin{proof}
By induction on $T$. The case $T=1$ is the hypothesis. Suppose
$R_T(x) = R_T(y)$. Since the dynamics are time-homogeneous,
$R_{T+1}(x) = \bigcup_{z \in R_T(x)} R_1(z)$ and likewise for $y$; since
$R_T(x) = R_T(y)$, these unions range over the same set of states $z$,
so $R_{T+1}(x) = R_{T+1}(y)$.
\end{proof}

\begin{remark}
This is the precise sense in which a collision, once it happens, is
permanent under time-homogeneous dynamics: two states whose immediate
futures coincide exactly will have every subsequent future coincide
as well, since there is nothing left in either state's future to
distinguish it from the other's.
\end{remark}

\section*{Quotient invariants}

\begin{proposition}[Quotient invariance]
\label{prop:invariance}
Let $F : \mathcal{K}(X) \to Y$ be any function. Then
$F \circ \Gamma : H \to Y$ factors through $q$: there is a unique
$\bar{F} : H/{\sim_\Gamma} \to Y$ with $F \circ \Gamma = \bar{F} \circ q$.
\end{proposition}
\begin{proof}
$F \circ \Gamma : H \to Y$ satisfies $(F \circ \Gamma)(h_1) = (F \circ \Gamma)(h_2)$
whenever $h_1 \sim_\Gamma h_2$, since $\Gamma(h_1) = \Gamma(h_2)$ by
definition of $\sim_\Gamma$ and $F$ is a function, hence single-valued.
Apply Proposition~\ref{prop:universal} with $F \circ \Gamma$ in place
of $\Gamma$ to obtain the factorization.
\end{proof}

\begin{corollary}
Continuation volume $\operatorname{Vol}_T$, the distinction margin
$A_{\mathcal D}^T$ relative to any fixed reference distinction, and
boundary proximity $\delta$, all defined in \emph{Continuation
Geometry} as functions of a reachable set or of a state's position
relative to one, descend to well-defined functions on
$H/{\sim_\Gamma}$.
\end{corollary}
\begin{proof}
Each is of the form $F(\Gamma(h))$ or $F(R_T(x(h)))$ for a function
$F$ depending only on the reachable set itself (Lebesgue measure of
the set, Hausdorff distance to another fixed set, ambient distance
from a fixed reference set), so Proposition~\ref{prop:invariance}
applies directly to each.
\end{proof}

\begin{remark}
This is a broader statement than it might first appear. It says that
any quantity this book or its predecessor has defined, or will define,
purely in terms of a reachable set, without reference to which
history produced that set, is automatically constant on
$\sim_\Gamma$-equivalence classes, with no further argument required
each time a new such quantity is introduced. The work of checking
invariance, done once here, does not need to be redone for every
future functional built the same way.
\end{remark}

\section*{Finite checkability}

Lemma~\ref{lem:propagate} shows that equal one-step images propagate
forward forever. It says nothing about the converse: two states whose
one-step images differ can still coincide at some later horizon.
Nothing in this chapter has claimed otherwise, and the general
question, whether continuation equivalence at every horizon can still
be checked in finitely many steps despite this, deserves a correct
answer rather than an implied one.

\begin{proposition}[Eventual periodicity bounds the check]
\label{prop:periodicity}
If $X$ is finite, then for any $x, y \in X$ the paired sequence
$\big(R_1(x),R_1(y)\big), \big(R_2(x),R_2(y)\big), \dots$ is eventually
periodic, with both the pre-period and the period bounded by
$4^{|X|}$. Consequently, whether $R_T(x) = R_T(y)$ holds for all
sufficiently large $T$ can be determined by comparing $R_T(x)$ and
$R_T(y)$ for $T$ ranging only up to $2 \cdot 4^{|X|}$.
\end{proposition}
\begin{proof}
Each pair $\big(R_T(x), R_T(y)\big)$ is an element of $2^X \times 2^X$,
which has $4^{|X|}$ elements, so the paired sequence takes at most
$4^{|X|}$ distinct values; among its first $4^{|X|}+1$ terms, two must
coincide by pigeonhole, $\big(R_p(x),R_p(y)\big) =
\big(R_q(x),R_q(y)\big)$ for some $1 \leq p < q \leq 4^{|X|}+1$.
Time-homogeneity means $R_{T+1} = \Psi(R_T)$ for the fixed map
$\Psi(S) := \bigcup_{z \in S} \Phi(z)$, applied identically to the $x$
and $y$ components, so equality of the pair at $p$ and $q$ propagates:
$\big(R_{p+j}(x), R_{p+j}(y)\big) = \big(\Psi^j(R_p(x)), \Psi^j(R_p(y))\big)
= \big(\Psi^j(R_q(x)), \Psi^j(R_q(y))\big) = \big(R_{q+j}(x), R_{q+j}(y)\big)$
for every $j \geq 0$: the \emph{paired} sequence is periodic with
period dividing $q - p \leq 4^{|X|}$ from index $p \leq 4^{|X|}$
onward. Comparing the pair over one further full period, at most
$4^{|X|}$ more steps beyond the pre-period bound, settles whether $x$
and $y$ agree from that point on.
\end{proof}

\begin{remark}
Bounding each sequence's own period separately and then comparing
them over a single shared window is not sufficient on its own: two
sequences with different individual periods $p_x, p_y$ need not
realign within $\max(p_x,p_y)$ steps, only within
$\operatorname{lcm}(p_x,p_y)$ steps, which can exceed either period
substantially. Working with the pair from the start avoids this,
since the pair's own period already accounts for however the two
component periods interact.
\end{remark}

\begin{remark}
This is the correct version of the claim that continuation
equivalence is finitely checkable. It is a coarser bound than a sharp
analysis of any particular system would give, since $2^{|X|}$ is a
worst case over all possible dynamics on a state space of that size,
but it is unconditional: no structural assumption on $\Phi$ beyond
finiteness of $X$ and time-homogeneity is needed. The traffic light
below reaches its own eventual behavior in four steps, far inside this
bound, and it is worth seeing exactly how.
\end{remark}

\section*{The traffic light, revisited}

Recall the traffic light of \emph{Continuation Geometry}, Chapter 10:
states $(k, r) \in \{0,1,2,3\} \times \{0,1\}$, phase $k$ and pending
request $r$. To compute $\Gamma$ properly requires treating the
system's branching explicitly rather than leaving it implicit. Let
$e \in \{0,1\}$ record whether a new pedestrian request arrives during
a given step (exogenous, not controlled by the phase). The transition
is
\[
\Phi((k,r), e) = \left( (k+1) \bmod 4,\ r' \right), \qquad
r' = \begin{cases}
0 & k+1 \equiv 0 \pmod 4 \text{ (clear on arrival at red)} \\
1 & k+1 \not\equiv 0 \pmod 4 \text{ and } (r = 1 \text{ or } e = 1) \\
0 & k+1 \not\equiv 0 \pmod 4 \text{ and } r = 0 \text{ and } e = 0.
\end{cases}
\]
Let $H$ be the set of finite paths of the system, and
$\Gamma(h) := R_1(x(h))$, the one-step reachable set from the
endpoint of $h$, ranging over $e \in \{0,1\}$.

\begin{proposition}[Exact computation of the traffic-light quotient at $T=1$]
\label{prop:tlquotient}
Under $\sim_\Gamma$, the eight states fall into exactly seven classes:
the single nontrivial class $\{(3,0),(3,1)\}$, and six singleton
classes, one for each remaining state.
\end{proposition}
\begin{proof}
Compute $R_1$ for all eight states directly from $\Phi$. For
$k \neq 3$, arrival is at a nonzero phase and the ordinary branching
rule applies, $r' = 1$ iff $r=1$ or $e=1$:
\begin{align*}
R_1(0,0) &= \{(1,0),(1,1)\}, & R_1(0,1) &= \{(1,1)\}, \\
R_1(1,0) &= \{(2,0),(2,1)\}, & R_1(1,1) &= \{(2,1)\}, \\
R_1(2,0) &= \{(3,0),(3,1)\}, & R_1(2,1) &= \{(3,1)\}.
\end{align*}
For $k=3$, arrival is at phase $0$, so clearing fires unconditionally
regardless of $r$ or $e$:
\[
R_1(3,0) = \{(0,0)\}, \qquad R_1(3,1) = \{(0,0)\}.
\]
Comparing all eight sets directly: the three size-two sets,
$\{(1,0),(1,1)\}$, $\{(2,0),(2,1)\}$, $\{(3,0),(3,1)\}$, are pairwise
distinct; the four remaining size-one sets, $\{(1,1)\}$, $\{(2,1)\}$,
$\{(0,0)\}$ (appearing twice, from $(3,0)$ and $(3,1)$), and
$\{(0,0)\}$ again, give exactly one coincidence, $R_1(3,0) = R_1(3,1)
= \{(0,0)\}$, and no other.
\end{proof}

Now define the mechanical quotient $\sim_M$ independently of
$\Gamma$, as the identification the design actually performs: the
clearing rule fires on the transition \emph{into} phase $0$, so it is
the two states about to make that transition, $(3,0)$ and $(3,1)$,
that the design treats as interchangeable, both being followed
unconditionally by $(0,0)$.

\begin{proposition}[The design is exact at $T=1$]
$\sim_M\ =\ \sim_\Gamma$.
\end{proposition}
\begin{proof}
Both relations have exactly one nontrivial class, $\{(3,0),(3,1)\}$,
and agree as singletons on every other state: $\sim_M$ by definition
of the clearing rule, $\sim_\Gamma$ by Proposition~\ref{prop:tlquotient}.
Two equivalence relations on the same finite set with identical
partitions are equal.
\end{proof}

\begin{remark}
It is worth being precise about which of two similar-sounding designs
this result actually vindicates. Clearing on arrival at phase $0$
(the rule above) is exact. A superficially similar alternative,
clearing whenever the system is currently \emph{at} phase $0$ rather
than arriving there, would instead identify $(0,0)$ and $(0,1)$. But
$R_1(0,0) = \{(1,0),(1,1)\} \neq \{(1,1)\} = R_1(0,1)$: a state at the
red light with no pending request can still branch toward either
outcome next, while one with a pending request cannot, precisely
because that request is about to be served. The alternative design
would merge two states with genuinely different one-step futures,
exactly the coarser, information-destroying case this chapter opened
by distinguishing from the exact one. The two designs differ only in
whether clearing is keyed to the transition or to the state, and only
one of them is correct.
\end{remark}

\section*{What one horizon hides}

Restricting Proposition~\ref{prop:tlquotient} to $T=1$ was not a
simplification; it is what $\Gamma$ was defined to measure. But
Proposition~\ref{prop:periodicity} guarantees the longer-horizon
question has a definite, finitely checkable answer, and for this
system that answer is worth computing directly rather than left as an
abstract guarantee.

\begin{proposition}[The request bit is eventually redundant everywhere]
\label{prop:eventual}
For each phase $k$, the pair $(k,0)$ and $(k,1)$ satisfies
$R_T(k,0) = R_T(k,1)$ for all $T \geq 4 - k$, and not before.
\end{proposition}
\begin{proof}[Proof by direct computation]
Compute $R_T(k,0)$ and $R_T(k,1)$ for $T = 0, \dots, 4-k$ at each
phase directly from $\Psi$, rather than inferring later horizons from
earlier ones: a size mismatch at one horizon does not, on its own,
preclude equality reappearing or vanishing at a different horizon, so
each phase's first point of equality is verified independently here
rather than propagated.

For $k=3$: $R_0(3,0) = \{(3,0)\} \neq \{(3,1)\} = R_0(3,1)$;
$R_1(3,0) = R_1(3,1) = \{(0,0)\}$ (Proposition~\ref{prop:tlquotient}).
First equality at $T=1=4-3$.

For $k=2$: $R_0(2,0) = \{(2,0)\} \neq \{(2,1)\} = R_0(2,1)$;
$R_1(2,0) = \{(3,0),(3,1)\} \neq \{(3,1)\} = R_1(2,1)$ (both sets
computed directly from $\Phi$, not equal since one has two elements
and the other one); $R_2(2,0) = \Psi(\{(3,0),(3,1)\}) = \{(0,0)\} =
\Psi(\{(3,1)\}) = R_2(2,1)$. First equality at $T=2=4-2$.

For $k=1$: $R_0, R_1$ inequality follows the same way (cardinality
$2$ versus $1$ at $T=1$, distinct singletons at $T=0$); $R_2(1,0) =
\{(3,0),(3,1)\} \neq \{(3,1)\} = R_2(1,1)$, still unequal;
$R_3(1,0) = \Psi(\{(3,0),(3,1)\}) = \{(0,0)\} = \Psi(\{(3,1)\}) =
R_3(1,1)$. First equality at $T=3=4-1$.

For $k=0$: $R_0(0,0) = \{(0,0)\} \neq \{(0,1)\} = R_0(0,1)$, distinct
singletons; $R_1, R_2, R_3$ each pit a two-element set against a
one-element set, as computed directly from $\Phi$ at each step, so
none of the three can coincide; $R_4(0,0) = \{(0,0)\} = R_4(0,1)$.
First equality at $T=4=4-0$.

In every case the inequality at each horizon strictly before $4-k$
was checked directly from the explicit sets, not inferred from the
inequality at a different horizon.
\end{proof}

\begin{remark}
This is the sharper and more interesting fact the single-horizon
result was hiding. The request bit is not redundant in general; it is
redundant \emph{by a deadline}, and that deadline is exactly the
number of steps remaining until the light is guaranteed to cycle
through its clearing transition, four steps for a state that has just
left it, one step for a state about to enter it. Every state's
request bit stops mattering to its continuation space at precisely
the moment continuation stops being able to tell whether that request
will still be pending when clearing finally arrives, and not one step
before. The design clears the bit at the earliest possible moment,
$T=1$ after arrival, which is also the latest moment at which failing
to clear it would still have been observably wrong; every other
phase's bit becomes redundant on its own, without any design
intervention, simply because the light's own periodicity eventually
answers the question the bit was tracking.
\end{remark}

\section*{What quotients do not yet explain}

This chapter has shown when two histories are interchangeable from
the standpoint of every quantity built purely from continuation
spaces, and it has shown, in one worked case, that a designed
identification can be checked against that standard exactly rather
than merely defended by intuition. It has not yet said anything about
what is lost when a quotient is coarser than $\sim_\Gamma$, only that
the traffic light's particular quotient happens not to be. The next
chapter takes up the coarser case directly: not whether information is
destroyed, which this chapter's machinery can already detect, but how
much, and by what measure.

\chapter{Distinction Loss Under Quotients}

A library catalog entry is a quotient of a book: title, author, and a
handful of subject headings standing in for everything between the
covers. A good catalog loses nothing a patron actually needs in order
to decide whether to pull the book from the shelf. A bad one can lose
quite a lot; a catalog that records only ``fiction'' or
``nonfiction'' will happily file a memoir next to a novel and a field
guide next to a textbook, and a patron relying on that heading alone
has no way to tell them apart until the book is already in hand. The
catalog entry is not wrong, exactly; it simply was not built to
preserve the distinctions this particular patron cares about.

The previous chapter showed how to check whether a given
identification of histories is exact: whether it merges only
histories whose continuation spaces already coincide, or whether it
merges some that do not. That check was binary, exact or not. This
chapter asks the sharper question a librarian actually has to answer.
When an identification is not exact, how much is lost, and where
exactly does the loss occur? The traffic light's badly designed
alternative from the previous chapter turns out to have a precise
answer, not just a qualitative one.

\section*{Fibers and their content}

\begin{definition}[Fiber, $\Gamma$-content, loss]
For a quotient $q : H \to O$ and $o \in O$, the fiber over $o$ is
$q^{-1}(o) \subseteq H$. Its $\Gamma$-content is
$C_\Gamma(o) := \{\Gamma(h) : h \in q^{-1}(o)\}$, the set of distinct
continuation spaces appearing in the fiber. The loss at $o$ is
\[
\mathcal{L}(o) := \log |C_\Gamma(o)|,
\]
so $\mathcal{L}(o) = 0$ exactly when every history mapping to $o$
under $q$ shares a single continuation space.
\end{definition}

\begin{proposition}[Zero loss characterizes $\Gamma$-respecting quotients]
\label{prop:zeroloss}
$\mathcal{L}(o) = 0$ for every $o \in O$ if and only if $q$ refines
$\sim_\Gamma$: $h_1 \sim_q h_2 \implies h_1 \sim_\Gamma h_2$.
\end{proposition}
\begin{proof}
If $q$ refines $\sim_\Gamma$, then $h_1 \sim_q h_2$ implies
$\Gamma(h_1) = \Gamma(h_2)$, so every fiber of $q$ lies inside a
single $\sim_\Gamma$-class and has $|C_\Gamma(o)| = 1$, giving
$\mathcal{L}(o) = 0$. Conversely, if $\mathcal{L}(o) = 0$ for every
$o$, then every fiber has $|C_\Gamma(o)| = 1$, so any two histories in
the same fiber share the same $\Gamma$-value: $h_1 \sim_q h_2 \implies
\Gamma(h_1) = \Gamma(h_2) \implies h_1 \sim_\Gamma h_2$, so $q$
refines $\sim_\Gamma$.
\end{proof}

\section*{The coarsest lossless quotient}

\begin{theorem}[$\sim_\Gamma$ is the coarsest zero-loss quotient]
\label{thm:coarsest}
Every zero-loss quotient $q : H \to O$ refines $\sim_\Gamma$, and
$\sim_\Gamma$ itself has zero loss everywhere. Consequently
$\sim_\Gamma$ is, up to relabeling of $O$, the unique coarsest
quotient of $H$ with zero loss.
\end{theorem}
\begin{proof}
The first claim is Proposition~\ref{prop:zeroloss}. That
$\sim_\Gamma$ itself has zero loss is immediate: by definition,
$h_1 \sim_\Gamma h_2$ already means $\Gamma(h_1) = \Gamma(h_2)$, so
every $\sim_\Gamma$-class has exactly one $\Gamma$-value. Since every
zero-loss $q$ refines $\sim_\Gamma$, no zero-loss quotient can be
strictly coarser than $\sim_\Gamma$, and $\sim_\Gamma$ achieves zero
loss, so it is the coarsest one.
\end{proof}

\begin{remark}
This is the quotient-theoretic form of a fact Chapter 5 of
\emph{Continuation Geometry} established for sufficient statistics
without this vocabulary. A sufficient statistic $\varphi$ there was
required to satisfy $\varphi(H_t) = \varphi(H_t') \implies
\Gamma(x_t \mid H_t) = \Gamma(x_t \mid H_t')$, which is exactly the
condition that the quotient induced by $\varphi$ refines
$\sim_\Gamma$. Theorem~\ref{thm:coarsest} says $\sim_\Gamma$ itself is
the coarsest possible sufficient statistic: the most compressed
summary of a history that still loses nothing continuation-relevant.
Any coarser summary is not merely less detailed; by
Proposition~\ref{prop:zeroloss}, it has strictly positive loss
somewhere.
\end{remark}

\begin{proposition}[Coarsening never decreases loss]
\label{prop:monotone-ch3}
If $q'$ refines $q$ (every $q'$-fiber is contained in some $q$-fiber),
then for every $q$-fiber $q^{-1}(o)$, expressed as a union of
$q'$-fibers $q'^{-1}(o'_1), \dots, q'^{-1}(o'_k)$,
\[
\mathcal{L}_q(o) \geq \max_i \mathcal{L}_{q'}(o'_i).
\]
\end{proposition}
\begin{proof}
$C_\Gamma(o) = \bigcup_i C_{\Gamma}(o'_i)$, since the $\Gamma$-values
appearing in $q^{-1}(o)$ are exactly those appearing in its
constituent $q'$-fibers. A union of sets is at least as large as any
one of them, so $|C_\Gamma(o)| \geq |C_\Gamma(o'_i)|$ for each $i$,
and $\log$ is monotone increasing.
\end{proof}

\begin{remark}
Merging two fibers can only add to the set of continuation spaces an
observer sees when looking at the merged fiber; it can never remove
one. This is the precise sense in which coarsening a quotient is a
one-way trip toward loss: refining a lossy quotient can repair the
loss, by splitting a mixed fiber back into pieces that each see only
one $\Gamma$-value, but no further coarsening of an already-lossy
quotient can undo the damage, only add to it or leave it unchanged.
\end{remark}

\section*{Distinction margin as a local loss detector}

\begin{proposition}[Vanishing margin detects a redundant distinction]
For a distinction $\mathcal{D}$ and $x \in K$, $A_{\mathcal D}^T(x) = 0$
(Chapter 7 of \emph{Continuation Geometry}) if and only if some $y$
outside $x$'s $\mathcal{D}$-cell satisfies $R_T(x) = R_T(y)$: that is,
$\mathcal{D}$ places $x$ and $y$ in different cells even though their
continuation spaces coincide at horizon $T$. Vanishing margin is a
diagnostic for $\mathcal{D}$ being \emph{finer} than $\sim_\Gamma$ at
$x$, drawing a distinction continuation geometry does not support, not
for $\mathcal{D}$ being coarser.
\end{proposition}
\begin{proof}
Recall $A_{\mathcal D}^T(x) := \inf_{y \notin D(x)} d_H(R_T(x), R_T(y))$.
This infimum is $0$ exactly when some $y$ outside $x$'s
$\mathcal{D}$-cell has $R_T(x) = R_T(y)$ (attained, under the
compactness hypothesis already used in Chapter 7), which is exactly
the statement that $x$ and that $y$ share a $\Gamma$-value at horizon
$T$ despite $\mathcal D$ having placed them in different cells: they
are $\sim_\Gamma$-equivalent even though $\mathcal D$ distinguishes
them. This is over-refinement, not collapse.
\end{proof}

\begin{definition}[Within-cell spread]
For a $\mathcal{D}$-cell $D$, $\operatorname{diam}_\Gamma(D) :=
\sup_{x,y \in D} d_H(R_T(x), R_T(y))$.
\end{definition}

\begin{proposition}[Zero spread is the actual collapse-safety detector]
\label{prop:diamzero}
$\operatorname{diam}_\Gamma(D) = 0$ if and only if $D$ is contained in
a single $\sim_\Gamma$-class at horizon $T$, i.e.\ $\mathcal{D}$
introduces no loss in the sense of the fiber-loss definition earlier
in this chapter, applied to $D$ as a fiber.
\end{proposition}
\begin{proof}
$\operatorname{diam}_\Gamma(D) = 0$ means $d_H(R_T(x),R_T(y)) = 0$ for
every $x,y \in D$, i.e.\ $R_T(x) = R_T(y)$ for every pair in $D$, i.e.\
$D$ lies inside one $\sim_\Gamma$-class; this is exactly
$|C_\Gamma(D)| = 1$ from the definition of fiber loss, i.e.\
$\mathcal{L}(D) = 0$.
\end{proof}


\begin{remark}
This gives the distinction margin of the first book a second reading,
and a corrective one. It was introduced there as a measure of how
close a state sits to losing a distinction, language that invites
exactly the confusion the propositions above are meant to head off.
$A_{\mathcal D}^T(x) = 0$ is not a report of quotient loss; it is a
report that $\mathcal{D}$ is drawing a line continuation geometry
does not, distinguishing two states with identical futures. Quotient
loss, in the sense the rest of this chapter measures it, is instead
the province of $\operatorname{diam}_\Gamma$: a cell $D$ with
$\operatorname{diam}_\Gamma(D) > 0$ is the one actually merging states
with different futures, the failure mode Chapter 3 of this book was
built to quantify. The two diagnostics point in opposite directions,
over-refinement and under-refinement, and a single distinction
$\mathcal{D}$ can exhibit either, both, or neither, independently at
different points of $K$.
\end{remark}

\section*{The traffic light, quantified}

Return to the two designs compared in the previous chapter: the
correct one, clearing on arrival at phase $0$, and the alternative,
clearing whenever currently at phase $0$.

\begin{example}[Zero loss, and one bit]
The correct design's quotient is $\sim_\Gamma$ itself (Proposition~4
of the previous chapter), so by Theorem~\ref{thm:coarsest} its loss is
$\mathcal{L}(o) = 0$ at every fiber $o$. The alternative design merges
$(0,0)$ and $(0,1)$ into a single fiber $o_{\mathrm{bad}}$. Their
$\Gamma$-values at $T=1$ are $R_1(0,0) = \{(1,0),(1,1)\}$ and
$R_1(0,1) = \{(1,1)\}$, which are distinct, so
$C_\Gamma(o_{\mathrm{bad}}) = \{\{(1,0),(1,1)\},\ \{(1,1)\}\}$, a set
of size $2$, giving
\[
\mathcal{L}(o_{\mathrm{bad}}) = \log 2,
\]
exactly one bit, and $\mathcal{L}(o) = 0$ at every other fiber, since
the alternative design agrees with the correct one everywhere except
at this single merge.
\end{example}

\begin{remark}
This is what the previous chapter's qualitative verdict, that the
alternative design was coarser and information-destroying, actually
amounts to once measured: not a diffuse degradation spread across the
whole system, but a single, exactly located bit of loss, sitting at
precisely the fiber where the alternative design merged two states
that a patron of this particular library, a controller trying to
plan four steps ahead, would still have wanted to tell apart. Every
other fiber of the alternative design is, by coincidence or by the
limited scope of its one design error, exactly as good as the correct
one.
\end{remark}

\section*{What loss does not yet explain}

This chapter can say, precisely, how much is lost by a given
quotient and exactly where. It has nothing to say about why a
particular coarser quotient might be chosen anyway, in full knowledge
of the loss it incurs. A designer might prefer a one-bit-lossy
traffic light if the saved bit of memory matters more than the
distinction it would have preserved; a librarian might prefer the
cheap two-category catalog if most patrons never need the finer one.
Measuring loss and deciding whether a loss is worth paying for are
different questions, and this chapter has only ever answered the
first of them.

\chapter{Continuation Operators}

The first three chapters of this book have all, in their own way, been
about nouns: a shape a system of contractions converges to, a
partition of histories that share a future, a number measuring what a
partition throws away. None of them named the verb. A contraction
does something to a shape. A design decision, clearing a bit on
arrival at a red light rather than while sitting at one, does
something to a history. This chapter is about the doing: what counts
as a legitimate move in this framework, and what survives when two
different moves are composed into one.

The two kinds of move already used in this book look, on the surface,
like opposites. Chapter 1's contractions built detail up, each
application adding a smaller, self-similar copy of the whole. Chapters
2 and 3's quotients tore detail down, collapsing histories that had
become, for every purpose this book has formalized, the same. It
would be a missed opportunity to leave them looking like two
unrelated tools. They are both instances of one idea, and this
chapter's job is to say what that idea is.

\section*{Operators and their composition}

\begin{definition}[Continuation operator]
A continuation operator on $X$ is a Lipschitz map
$T : \mathcal{K}(X) \to \mathcal{K}(X)$: there is a constant
$L \geq 0$ with $d_H(T(A), T(B)) \leq L\, d_H(A,B)$ for all
$A, B \in \mathcal{K}(X)$. Write $L(T)$ for the smallest such constant.
\end{definition}

\begin{proposition}[Closure under composition]
\label{prop:closure}
If $T_1, T_2$ are continuation operators, so is $T_2 \circ T_1$, with
$L(T_2 \circ T_1) \leq L(T_2)\, L(T_1)$.
\end{proposition}
\begin{proof}
$d_H(T_2(T_1(A)), T_2(T_1(B))) \leq L(T_2)\, d_H(T_1(A), T_1(B)) \leq
L(T_2)\, L(T_1)\, d_H(A,B)$.
\end{proof}

\begin{remark}
This proof is one line because it has to be; the definition was
chosen precisely so that composition would cost nothing further to
verify. What makes the definition worth having is not this
proposition but what it is a common home for. The contractions of
Chapter 1 are continuation operators with $L(T) < 1$. The map
$\Psi(A) := \bigcup_{z \in A} \Phi(z)$ used throughout Chapters 2 and
3 to advance a reachable set by one step is a continuation operator
too, generally with $L(\Psi) \geq 1$: it need not contract, but it is
still Lipschitz whenever the underlying one-step map $\Phi$ is well
behaved. Both belong to the same category. What distinguishes a
generative operator from a reductive one is not the definition above;
it is what happens under repeated composition, contraction toward a
single attractor in one case, and, as the next section makes precise,
collapse of distinctions in the other.
\end{remark}

\section*{Operators on histories, and what they respect}

The operators above act on shapes directly. A different and equally
natural kind of operator acts on histories: appending one more real
transition, editing a codebase, applying a repair. The interesting
question for an operator of this kind is not its Lipschitz constant
but whether it interacts predictably with $\sim_\Gamma$.

\begin{definition}[$\Gamma$-compatible operator]
$T : H \to H$ is $\Gamma$-compatible if there is a map
$\hat{T} : \mathcal{K}(X) \to \mathcal{K}(X)$ with
$\Gamma(T(h)) = \hat{T}(\Gamma(h))$ for every $h \in H$: the
continuation space of the transformed history depends on the original
history only through its continuation space, not otherwise.
\end{definition}

\begin{proposition}[$\Gamma$-compatible operators descend to the quotient]
\label{prop:descend}
If $T$ is $\Gamma$-compatible, then $h_1 \sim_\Gamma h_2 \implies
T(h_1) \sim_\Gamma T(h_2)$: $T$ induces a well-defined map
$\bar{T} : H/{\sim_\Gamma} \to H/{\sim_\Gamma}$.
\end{proposition}
\begin{proof}
This is Proposition~4 of the previous chapter (quotient invariance),
applied to $F := \hat{T} \circ \Gamma$ in place of $\Gamma$ itself:
$h_1 \sim_\Gamma h_2$ gives $\Gamma(h_1) = \Gamma(h_2)$, hence
$\Gamma(T(h_1)) = \hat{T}(\Gamma(h_1)) = \hat{T}(\Gamma(h_2)) =
\Gamma(T(h_2))$, so $T(h_1) \sim_\Gamma T(h_2)$.
\end{proof}

\begin{remark}
This is the sense in which $\Gamma$-compatibility is the right
condition to ask for: it says an operator cannot secretly depend on
information that continuation equivalence has already declared
irrelevant. An operator failing this condition would be able to treat
two continuation-equivalent histories differently for reasons that
have nothing to do with what either history can still become, which
is exactly the kind of dependence Chapter 5 of \emph{Continuation
Geometry} worried about under the name of history-dependence in the
first place.
\end{remark}

\begin{example}[The traffic light's step operators]
For $e \in \{0,1\}$, let $T_e : H \to H$ append the transition
$\Phi(\cdot, e)$ to a history. Whether $T_e$ is $\Gamma$-compatible
(at horizon $T=1$) needs checking only on the one nontrivial
$\sim_\Gamma$-class found in Chapter 2, $\{(3,0),(3,1)\}$, since every
other class is a singleton and compatibility there is automatic.
Direct computation gives $T_0(3,0) = T_0(3,1) = (0,0)$ and
$T_1(3,0) = T_1(3,1) = (0,0)$: both step operators send the two
equivalent histories to the identical resulting state, not merely to
states with equal $R_1$, so compatibility holds here in the strongest
possible form.
\end{example}

\begin{remark}
This is a narrow check, confined to the single class where there was
anything to check, and it should not be read as a general fact about
$T_0$ or $T_1$ on this system, only as a confirmation that neither
step operator introduces a new dependence on which of $(3,0)$ or
$(3,1)$ a history happened to pass through, at exactly the point
where such a dependence would have mattered.
\end{remark}

\section*{The family $\Gamma_T$ and the operator $\Psi$}

A different relationship, easy to conflate with $\Gamma$-compatibility
but distinct from it, has been used without a name since Chapter 2. It
concerns not an operator on histories but the relationship between the
finite-horizon continuation spaces at consecutive horizons for one and
the same history.

\begin{remark}[$\Psi$ relates horizons, not histories]
By construction, $R_{T+1}(x) = \bigcup_{z \in R_T(x)} \Phi(z) =
\Psi(R_T(x))$, i.e.\ $\Gamma_{T+1} = \Psi \circ \Gamma_T$ as an
identity between functions $H \to \mathcal{K}(X)$, where
$\Gamma_T(h) := R_T(x(h))$. This is not an instance of
Proposition~\ref{prop:descend}; no operator on $H$ is involved, only
the fixed continuation operator $\Psi$ acting on the shape side, once
for each increment of the horizon. It is this identity, not any
property of an operator on histories, that did the real work in
Lemma~2 of the previous chapter (one-step agreement propagates) and in
the eventual-periodicity argument of Chapter 3: both are, in
retrospect, statements about the orbit of a single reachable set
under repeated application of the one continuation operator $\Psi$,
dressed in the language of horizons rather than the language of
operators used here.
\end{remark}

\section*{Two families, one category}

\begin{corollary}
Both the contractions of Chapter 1 and the operator $\Psi$ of Chapters
2 and 3 are continuation operators in the sense of this chapter's
first definition; they differ only in whether their Lipschitz
constant is below $1$.
\end{corollary}

\begin{remark}
It is tempting to read this as saying generation and collapse are two
regions of one parameter, contraction below $1$, something else
above it, and that reading is too strong. $L(T) < 1$ does imply
convergence to a unique attractor, by Chapter 1's own theorem; that
direction is solid. $L(T) \geq 1$ implies nothing in particular. The
identity map has $L=1$ and collapses nothing. An isometry has $L=1$
and collapses nothing. A map that stretches every distance by a
factor of two has $L=2$ and, again, collapses nothing, since
stretching does not identify points that were previously distinct.
What actually collapses distinctions is a property orthogonal to the
Lipschitz constant: whether the operator is injective, whether
distinct inputs can share an output. $\Psi$ collapses states in the
traffic light not because $L(\Psi) \geq 1$, but because the specific
map $\Phi$ underlying it is non-injective at the clearing transition,
sending both $(3,0)$ and $(3,1)$ to $(0,0)$; a different one-step map
with $L(\Psi) \geq 1$ but no collisions anywhere would generate no
quotient at all. The Lipschitz constant and injectivity are
independent properties an operator can carry in any combination, and
conflating them is exactly the kind of overstatement this book has
tried, elsewhere, to catch in itself before it went to print.
\end{remark}

\begin{remark}[A more honest organizing geometry]
What Chapters 1 through 3 actually populate is not a line but a small
space of independent properties: metric sensitivity ($L(T)$, whether
and how strongly the operator contracts or expands distances);
injectivity (whether the operator identifies distinct states,
independent of $L(T)$); and compatibility with $\Gamma$ (Chapter 4's
own $\Gamma$-compatibility condition, a further independent property
an operator either has or lacks). A contraction with $L<1$ and full
injectivity, the pure case of Chapter 1, generates a self-similar
attractor without collapsing anything along the way. A non-injective
$\Psi$ with $L(\Psi) \geq 1$, the traffic light's case, collapses
distinctions and does not converge to a point, cycling instead. Other
combinations, injective maps with $L \geq 1$, non-injective
contractions, are equally coherent and simply were not this book's
running examples. The genuine unification this chapter earned is that
all of these live in one vocabulary, continuation operators, checkable
by the same closure proposition regardless of where in this space
they sit. The overclaim to retract is that a single number tells you
where in that space an operator is.
\end{remark}

\section*{What operators do not yet explain}

This chapter has said what a legitimate move is and shown that both
of this book's central constructions are instances of it, but it has
not asked which finite collections of operators are rich enough to
produce every continuation geometry worth having, only that individual
operators compose predictably once they exist. That question, closer
to the universal-generation intuition this project began with, needs
the vocabulary built here as a prerequisite, and is not attempted in
this chapter.

\chapter{Toward Universal Generation}

A handful of letters generates every word in a dictionary that has
not been written yet. A handful of chess rules generates every game
that will ever be played under them, including games no one has
conceived of. A small, fixed vocabulary reaching an unbounded space
of outcomes is one of the oldest tricks available to any generative
system, and it is worth asking, now that this book has a settled
vocabulary for what a continuation operator is, whether continuation
geometries have the same property. Can a small, fixed family of
operators, chosen once and never enlarged, get arbitrarily close to
any admissible continuation geometry at all?

This chapter does not answer that question. It draws its boundary
precisely enough to say exactly what would need to be true for the
answer to be yes, proves the one direction that is actually within
reach, and is honest about the direction that is not. The gap between
those two is real mathematics, not a gap in effort, and closing it is
left where it belongs: open.

\section*{The conjecture, stated precisely}

\begin{conjecture}[Universal continuation generation]
\label{conj:universal}
There exists a finite family of point contractions $T_1, \dots, T_n :
X \to X$ and a single seed set $P \in \mathcal{K}(X)$, both fixed once
and independent of any particular target, such that for every
nonempty compact $A \subseteq X$ and every $\varepsilon > 0$, some
finite word $i_1, \dots, i_k$ satisfies
\[
d_H\big(A,\ (T_{i_k} \circ \cdots \circ T_{i_1})(P)\big) < \varepsilon,
\]
where $(T_{i_k} \circ \cdots \circ T_{i_1})(P) := \{(T_{i_k} \circ
\cdots \circ T_{i_1})(p) : p \in P\}$ is the composed point-map,
itself a single contraction, applied setwise to the fixed seed.
\end{conjecture}

\begin{remark}[Why the seed must be fixed]
Allowing the starting set to vary with $A$ would make the conjecture
trivial: given any target $A$ and any $\varepsilon$, take the starting
set to be an $\varepsilon$-net of $A$ itself, apply the identity word,
and the approximation is immediate regardless of what the $T_i$ are.
The entire content of the conjecture is that one seed and one
generator family, chosen before any target is known, must work for
every target afterward; fixing $P$ once, alongside the $T_i$, is not
an optional refinement of the statement, it is the statement.
\end{remark}

\begin{remark}
This is stated as a conjecture and not as a theorem because it is
one; nothing in this book has proved it, and nothing in this chapter
will. What can be made precise is why it is hard, and that requires
first showing what is \emph{not} hard.
\end{remark}

\section*{The Collage Theorem}

\begin{theorem}[Collage Theorem]
\label{thm:collage}
Let $T_1, \dots, T_n$ be contractions with $\mathcal{H}(A) =
\bigcup_i T_i(A)$ and maximum contraction constant $c < 1$, and let
$A^\ast$ be the attractor of $\mathcal{H}$. Then for any nonempty
compact $A$,
\[
d_H(A, A^\ast) \leq \frac{d_H(A, \mathcal{H}(A))}{1-c}.
\]
\end{theorem}
\begin{proof}
By the triangle inequality and $A^\ast = \mathcal{H}(A^\ast)$,
\[
d_H(A, A^\ast) \leq d_H(A, \mathcal{H}(A)) + d_H(\mathcal{H}(A), \mathcal{H}(A^\ast)).
\]
By Proposition~3 of Chapter 1 ($\mathcal{H}$ is a contraction on
shapes), $d_H(\mathcal{H}(A), \mathcal{H}(A^\ast)) \leq c\, d_H(A,
A^\ast)$. Substituting,
\[
d_H(A, A^\ast) \leq d_H(A, \mathcal{H}(A)) + c\, d_H(A, A^\ast),
\]
and rearranging, since $c < 1$, gives
$(1-c)\, d_H(A,A^\ast) \leq d_H(A, \mathcal{H}(A))$, which is the
claim.
\end{proof}

\begin{remark}
Nothing here was invented for this chapter; the theorem is a direct
three-line consequence of facts Chapter 1 already established, and it
is worth noting how much it buys for how little it costs. It converts
a question about the attractor, an object defined only as a limit, so
generally hard to reason about directly, into a question about a
single application of $\mathcal{H}$ to a set already in hand. Finding
an operator family whose attractor is close to a target $A$ reduces to
finding one whose \emph{one-step collage} $\mathcal{H}(A)$ is already
close to $A$, an ordinary geometric fitting problem rather than a
limiting one.
\end{remark}

\section*{Existence of an approximating family}

\begin{theorem}[Every compact target is approximable]
\label{thm:existence}
For every nonempty compact $A \subseteq X = \mathbb{R}^n$ and every
$\varepsilon > 0$, there exists a finite family of contractions
(depending on $A$ and $\varepsilon$) whose attractor $A^\ast$
satisfies $d_H(A, A^\ast) < \varepsilon$.
\end{theorem}
\begin{proof}
Let $\{x_1, \dots, x_m\} \subseteq A$ be a finite $\delta$-net of $A$
with $\delta < \varepsilon$, which exists since $A$ is compact, hence
totally bounded, and can be chosen with net points drawn from $A$
itself. For each $i$, let $T_i : X \to X$ be the constant map
$T_i(x) := x_i$. Each $T_i$ is trivially a contraction, with constant
$c = 0$, since $d(T_i(x), T_i(y)) = 0 \leq 0 \cdot d(x,y)$ for all
$x,y$. Then $T_i(A) = \{x_i\}$ exactly, for any nonempty $A$, so
$\mathcal{H}(A) = \bigcup_i T_i(A) = \{x_1, \dots, x_m\}$, the net
itself. Since the net points lie in $A$ and every point of $A$ is
within $\delta$ of the net, $d_H(A, \mathcal{H}(A)) = d_H(A,
\{x_1,\dots,x_m\}) \leq \delta$. By Theorem~\ref{thm:collage} with
$c=0$, $d_H(A, A^\ast) \leq d_H(A,\mathcal{H}(A))/(1-0) \leq \delta <
\varepsilon$.
\end{proof}

\begin{remark}
This proof is fully elementary once constant maps are admitted as
contractions, which they trivially are under the definition used
throughout this book, and it is worth saying so plainly rather than
gesturing at a more elaborate construction that turns out not to be
needed. A version restricted to contractions with ratio bounded away
from $0$ demands more care, positioning each $T_i$ so its image lands
near $x_i$ without collapsing to a point, and is the version
classically used in fractal image compression \cite{barnsley1988};
that version is not reproduced here, since the elementary one already
proves everything this theorem claims.
\end{remark}

\begin{remark}
This is real progress, and it is worth being precise about exactly
how much. It says the space of Hutchinson attractors is dense, in
Hausdorff distance, in the space of nonempty compact subsets of $X$:
nothing worth generating is permanently out of reach. What it does
not say, and what Conjecture~\ref{conj:universal} actually asks for,
is that a \emph{single} family, fixed before any target is known,
suffices for every target. The construction above produces a
different family for each new target and each new tolerance; as
$\varepsilon \to 0$ or as $A$ grows more intricate, the number of
generators $m$ used in the proof generally grows without bound. The
existence theorem answers ``can this be reached at all,'' which
Theorem~\ref{thm:existence} settles completely. The conjecture asks
``can it be reached with tools chosen in advance,'' which is a
categorically stronger demand, and this chapter has not moved that
needle.
\end{remark}

\section*{What kind of gap this is}

The distinction is not a technicality, and it has a name in
neighboring fields worth being explicit about rather than gestured at.
The Stone--Weierstrass theorem \cite{stone1948} guarantees that polynomials are dense
in the continuous functions on a compact interval: a fixed class of
building blocks, closed under the operations that generate it, dense
in everything worth approximating. A universal approximation theorem
for neural networks \cite{cybenko1989} says something structurally similar: networks
built from a fixed activation function, with enough width, are dense
in the continuous functions. Both are theorems about a fixed generator
class being rich enough for every target, which is exactly the shape
of Conjecture~\ref{conj:universal} and exactly what
Theorem~\ref{thm:existence} does not provide, since its generator
family is chosen after the target is known.

\begin{remark}
Whether an analogue of Stone--Weierstrass or universal approximation
holds for continuation geometries specifically is not something this
chapter's methods can settle, and it should not be assumed to transfer
just because the surface shape of the claim resembles those theorems.
Both classical results rely on structure, algebraic closure under
multiplication for Stone--Weierstrass, compositional depth for neural
networks, that has no established analogue here. Building that
structure, or showing it cannot exist, is what would actually resolve
Conjecture~\ref{conj:universal}, and it remains the open question this
book set out to sharpen rather than close.
\end{remark}

\chapter{State Insufficiency, Made Explicit}

A doctor at the bedside observes breathing, heart rate, pupil
response, the ordinary vital signs a chart records every hour. Two
patients can match on every one of them, one sleeping, one comatose,
and the chart alone will not say which. This is not a failure of the
instruments. It is a fact about what the instruments were built to
see. The chart records an observation, a coarse summary of a much
richer state, and coarse summaries are allowed to erase exactly the
distinctions that happen not to show up in whatever they were designed
to measure. Whether a given summary erases a distinction that
matters, mere resemblance in the readings or a genuine difference in
what each patient can still become, is not a question the chart can
answer about itself.

This idea has been the motivating image behind this project since
before either book in it had a name, and it has never been made into
a theorem. It does not need a new toy system to become one. The
traffic light, already built and already verified twice over in
Chapters 2 and 3, contains a clean instance of exactly this
phenomenon, once it is looked at through a coarser observation than
the one this book has been using.

\section*{Observations and sufficiency}

\begin{definition}[Observation map, state sufficiency]
An observation map is any $\Pi : X \to O$. $\Pi$ is state-sufficient
(for continuation, at horizon $T$) if $\Pi(x) = \Pi(y) \implies
R_T(x) = R_T(y)$ for all $x,y \in X$; otherwise $\Pi$ is
state-insufficient.
\end{definition}

\begin{remark}
This is deliberately the same shape of question Chapter 3 already
answered in general. $\Pi$ is a quotient of $X$ in exactly the sense
used throughout this book, and state sufficiency is exactly zero loss
in the sense of that chapter, transported from quotients of histories
to observation maps on states. The next proposition makes the
translation exact rather than merely suggestive.
\end{remark}

\begin{proposition}[Sufficiency is refinement of $\sim_\Gamma$]
$\Pi$ is state-sufficient at horizon $T$ if and only if $\Pi$ refines
$\sim_\Gamma$ at that horizon: $\Pi(x) = \Pi(y) \implies x \sim_\Gamma
y$.
\end{proposition}
\begin{proof}
Immediate from the definitions: $x \sim_\Gamma y$ at horizon $T$ means
$R_T(x) = R_T(y)$, which is exactly the consequent of state
sufficiency's defining implication.
\end{proof}

\section*{An explicit failure of sufficiency}

\begin{definition}
Let $\Pi : \{0,1,2,3\} \times \{0,1\} \to \{0,1,2,3\}$,
$\Pi(k,r) := k$, the observation that records the traffic light's
phase but not its pending-request bit.
\end{definition}

\begin{theorem}[State insufficiency, explicit instance]
\label{thm:insufficient}
$\Pi$ is state-insufficient at $T=1$: $\Pi(0,0) = \Pi(0,1) = 0$, but
$R_1(0,0) = \{(1,0),(1,1)\} \neq \{(1,1)\} = R_1(0,1)$.
\end{theorem}
\begin{proof}
Both values of $\Pi$ are immediate from the definition of $\Pi$. The
inequality of the reachable sets is the computation already carried
out in Chapter 2 (Proposition 3, the $T=1$ traffic-light computation):
$R_1(0,0)$ and $R_1(0,1)$ were shown there to differ, as one of the
seven singleton classes and part of the argument establishing that no
merge occurs at phase $0$.
\end{proof}

\begin{remark}
This is the sleep/coma image with every part of it named. $\Pi$ plays
the role of the chart: it records the phase, the part of the state a
casual glance would call ``where things currently stand,'' and
discards the request bit, the part recording something about how the
state got there and what it is still owed. Two states indistinguishable
under $\Pi$, $(0,0)$ and $(0,1)$, are not indistinguishable under
$\Gamma$: one of them can still become $(1,0)$, a future the other
cannot reach at all. The chart is not lying. It was simply never built
to record the fact that would have mattered here.
\end{remark}

\section*{How much sufficiency fails, and where}

Chapter 3 built exactly the machinery needed to go beyond a single
witnessing pair and measure the full extent of $\Pi$'s insufficiency.

\begin{proposition}[The loss of $\Pi$, computed exactly]
For each phase $k \in \{0,1,2,3\}$, the loss at fiber $\Pi^{-1}(k) =
\{(k,0),(k,1)\}$ is
\[
\mathcal{L}(\Pi^{-1}(k)) = \begin{cases} \log 2 & k \in \{0,1,2\} \\ 0 & k = 3. \end{cases}
\]
\end{proposition}
\begin{proof}
By definition, $\mathcal{L}(\Pi^{-1}(k)) = \log|\{R_1(k,0), R_1(k,1)\}|$.
From the values already established (Chapter 2, Proposition 3, and
Theorem~\ref{thm:insufficient} above): $R_1(0,0) \neq R_1(0,1)$,
$R_1(1,0) = \{(2,0),(2,1)\} \neq \{(2,1)\} = R_1(1,1)$, and
$R_1(2,0) = \{(3,0),(3,1)\} \neq \{(3,1)\} = R_1(2,1)$, giving
$|\{\cdot,\cdot\}| = 2$ and loss $\log 2$ in each case. At $k=3$,
$R_1(3,0) = R_1(3,1) = \{(0,0)\}$, giving $|\{\cdot\}| = 1$ and loss
$0$.
\end{proof}

\begin{remark}
The pattern is worth reading rather than only computing. The
observation $\Pi$ loses exactly one bit at every phase except the one
immediately before clearing, where it loses nothing, for exactly the
reason Chapter 3's eventual-redundancy result already identified: at
phase $3$, the request bit is one step from becoming irrelevant
regardless of what $\Pi$ does, so discarding it early costs nothing
there and costs a full bit everywhere else. State insufficiency, once
quantified rather than merely demonstrated, is not a uniform fog
sitting over the whole system; it has a precise shape, concentrated
exactly where the discarded information still had work left to do.
\end{remark}

\section*{What this does and does not establish about sleep and coma}

Nothing above is a claim about physiology. It is a claim about one
eight-state machine and one particular coarsening of it, chosen
because the numbers needed to state it precisely were already on hand
and already checked. Whether any actual clinical observation, a chart,
a scan, a bedside exam, is state-insufficient for the specific
continuation spaces that distinguish sleep from coma is an empirical
question about biology, not something a toy traffic light can settle
by analogy. What this chapter has shown is narrower and more useful
than a claim about medicine: that state insufficiency is a real,
precisely statable, and exactly measurable property of an observation
map relative to a dynamics, that it can fail in some places and not
others within the very same system, and that the machinery already
built in this book, with no further construction, is already enough
to find out which.

\chapter{Invariant Decomposition, Staleness, and Intervention}

A scouting report can be wrong in two entirely different ways. It can
be out of date: accurate when filed, and steadily less trustworthy
with every day that passes since, repairable by sending another scout.
Or it can be asking the wrong question entirely: a report on one
trading region says nothing, however fresh, about a second region no
ship has ever crossed into, because the two have never exchanged a
single price, a single rumor, a single unit of anything. No amount of
resurveying the first region brings the second any closer to being
known. These look like the same complaint, an observer failing to
keep up with reality, and they are not. One is a fact about elapsed
time. The other is a fact about the shape of the space itself. This
chapter gives each its own formal treatment, because giving them the
same one would hide the fact that they need different remedies.

\section*{Invariant decomposition}

\begin{definition}[Invariant block]
$A \subseteq X$ is an invariant block under $F$ if $R_1(x) \subseteq
A$ for every $x \in A$.
\end{definition}

\begin{proposition}[Invariance propagates]
\label{prop:invpropagate}
If $A$ is an invariant block, then $R_T(x) \subseteq A$ for every $x
\in A$ and every $T \geq 0$.
\end{proposition}
\begin{proof}
By induction on $T$. $T=0$ is immediate ($x \in A$). If $R_T(x)
\subseteq A$, then $R_{T+1}(x) = \bigcup_{z \in R_T(x)} R_1(z)
\subseteq \bigcup_{z \in A} R_1(z) \subseteq A$, the last inclusion
by the definition of an invariant block applied to each $z \in A$.
\end{proof}

\begin{definition}[Invariant decomposition]
$X$ admits an invariant decomposition into $\{X_i\}$ if the $X_i$
partition $X$ and each $X_i$ is an invariant block.
\end{definition}

\begin{corollary}[Continuation confinement]
\label{cor:confine}
If $x \in X_i$, then $\Gamma_K(x) \subseteq \{\gamma : \gamma(t) \in
X_i \ \forall t\}$: every admissible continuation of $x$ remains in
$X_i$ forever.
\end{corollary}
\begin{proof}
Immediate from Proposition~\ref{prop:invpropagate}: $R_T(x) \subseteq
X_i$ for every $T$, and a trajectory whose position at every finite
horizon lies in $X_i$ lies in $X_i$ at every time.
\end{proof}

\begin{remark}
This is the precise sense in which observation in one block cannot
substitute for observation in another. For $x \in X_1$, $\Gamma_K(x)$
is determined entirely by the dynamics restricted to $X_1$; nothing
about $X_2$, however thoroughly sampled, enters into $R_T(x)$ for any
$T$, because Corollary~\ref{cor:confine} rules out any trajectory from
$x$ ever visiting $X_2$ in the first place. This is not a claim that
$X_2$ is hard to observe. It is a claim that $X_2$ is, formally,
irrelevant to $x$'s continuation space once the decomposition holds,
in the same sense that a variable a function does not depend on is
irrelevant to computing it, however precisely that variable is
measured. The classical distinction between a viability kernel
(admitting \emph{some} trajectory that stays put) and an invariance
kernel (requiring \emph{every} trajectory to stay put), already
present in the literature this book has cited since Chapter 1
\cite{aubinframkowska1990}, is exactly the distinction an invariant
block relies on, and exactly why this section is named for invariance
rather than for ergodicity: nothing here involves a measure or an
indecomposability condition on invariant measures, the actual content
of the word ergodic in its established sense. What is built here is
the purely topological half of that classical picture.
\end{remark}

\section*{Staleness}

Invariant decomposition is a permanent, structural fact about a system.
Staleness is not: it is what happens to an observer's knowledge of a
single, well-connected system as time passes without a fresh look.

\begin{definition}[Staleness bound]
Suppose $\dot x = f(x, u(t))$ for a fixed perturbation space $Q$
(independent of $x$) and $u(t) \in Q$ measurable, with $f$ Lipschitz
in $x$ with constant $L$ uniformly in $u \in Q$:
$\|f(x,u) - f(y,u)\| \leq L\|x-y\|$ for all $u \in Q$. Given an
observation fiber $\Pi^{-1}(o)$ of diameter $d_0 := \sup_{x,y \in
\Pi^{-1}(o)} d(x,y)$, the staleness bound at elapsed time $\Delta t$
is $d_0\, e^{L \Delta t}$.
\end{definition}

\begin{proposition}[Staleness grows at most exponentially]
\label{prop:staleness}
For $x, y \in \Pi^{-1}(o)$ and trajectories $\gamma_x, \gamma_y$
solving $\dot\gamma_x = f(\gamma_x, u(t))$, $\dot\gamma_y = f(\gamma_y,
u(t))$ from $x, y$ under the same measurable input $u(t) \in Q$,
$d(\gamma_x(\Delta t), \gamma_y(\Delta t)) \leq d_0\, e^{L \Delta t}$.
\end{proposition}
\begin{proof}
Let $w(t) := d(\gamma_x(t), \gamma_y(t))$. Since $f$ is $L$-Lipschitz
in its first argument uniformly in $u \in Q$ and both trajectories
are driven by the same $u(t)$,
$\dot{w}(t) \leq \|f(\gamma_x(t),u(t)) - f(\gamma_y(t),u(t))\| \leq
L\, w(t)$, so $w(t) \leq w(0) + L \int_0^t w(s)\, ds$. By the
continuous Gr\"onwall inequality (Appendix C of \emph{Continuation
Geometry}), $w(t) \leq w(0)\, e^{Lt} = d_0\, e^{Lt}$, giving the claim
at $t = \Delta t$.
\end{proof}

\begin{remark}
This is stated for single-valued dynamics with a shared perturbation
space, deliberately narrower than the set-valued $F$ used elsewhere in
this book. A genuine differential inclusion, where the two
trajectories are not assumed to realize a common measurable selection,
needs more than a Hausdorff--Lipschitz bound on $F$ to license the
same comparison estimate; it needs an explicit coupling, typically
supplied by a selection theorem of the kind Appendix B of
\emph{Continuation Geometry} already cites, or by a viability-theoretic
argument constructing matched trajectories directly. That extension is
not carried out here.
\end{remark}

\begin{remark}
Nothing in this proof is new; it is the same Gr\"onwall estimate
already used to bound trajectory growth under linear growth
hypotheses, applied here to the \emph{difference} between two
trajectories instead of the norm of one. What makes it worth restating
in this form is the reading it supports: staleness is not a fixed
penalty for not having observed recently, it is a bound that can
always be reset. Taking a fresh observation at time $\Delta t$
collapses the relevant fiber back down to whatever its new diameter
is, and the exponential clock restarts from there. This is the formal
content behind ``staleness is repaired by resampling'': the bound in
Proposition~\ref{prop:staleness} is a function of elapsed time since
the \emph{last} observation, not of any permanent property of the
system.
\end{remark}

\begin{remark}[Why the two do not reduce to one another]
A system separated into invariant blocks can be resampled as often as an
observer likes without ever closing the gap between blocks, since no
finite or infinite sequence of observations confined to $X_1$ changes
$\Gamma_K(x)$ for $x \in X_1$, by Corollary~\ref{cor:confine}. A stale
observation of a well-connected system closes its own gap with a
single fresh sample. Treating both under one heading, representation
degrading relative to territory, obscures that one is fixed by waiting
and looking again and the other cannot be fixed by looking again at
all, only by being, or gaining access to, the other block.
\end{remark}

\section*{Projection and intervention}

Chapter 6 formalized when an observation $\Pi$ fails to determine a
state's continuation space. It said nothing about what, if anything,
is done to a state once observed. The two are logically separate
questions.

\begin{definition}[Intervention operator]
An operator $T$ (Chapter 4) is an intervention on a set $S \subseteq
X$ if there is a fixed target $A_0 \in \mathcal{K}(X)$ with $\Gamma(T(x))
= A_0$ for every $x \in S$, regardless of $\Gamma(x)$.
\end{definition}

\begin{proposition}[Sufficiency and intervention are independent]
There exist systems in which $\Pi$ is state-sufficient on $S$ yet an
intervention operator is applied on $S$ anyway, and systems in which
$\Pi$ is state-insufficient on $S$ with no intervention applied at
all.
\end{proposition}
\begin{proof}
For the first: take $\Pi$ to be the identity (trivially sufficient,
since $\Pi(x)=\Pi(y) \implies x=y \implies \Gamma(x)=\Gamma(y)$), and
let $T$ map every $x \in S$ to a fixed state $x_0$ regardless of $x$;
$T$ is an intervention by construction, and $\Pi$'s sufficiency is
unaffected by whether $T$ is applied, since sufficiency is a property
of $\Pi$ alone. For the second: the traffic light's phase-only $\Pi$
of Chapter 6 is state-insufficient on $\{(0,0),(0,1)\}$, and the
system's actual dynamics $\Phi$ is not an intervention in the sense
above, since $\Gamma(\Phi(0,0)) \neq \Gamma(\Phi(0,1))$ in general
(the two states are not forced to a common target).
\end{proof}

\begin{remark}
This is a narrow, formal point, and it is worth being honest about how
narrow: it shows the two properties are not logically entailed by one
another, not that they are uncorrelated in how real institutions
actually behave. An institution whose observation is coarse may find
uniform treatment easier to justify or administer than
individually-tailored treatment would be, which is a pressure toward
intervention that has nothing to do with logical necessity and
everything to do with what coarse observation makes practical. The
independence proposition above rules out the claim that better
observation \emph{requires} less intervention or that worse
observation \emph{requires} more; it says nothing about which
correlations actually hold in systems built by people operating under
those practical pressures, which is a separate, empirical question
this chapter's methods do not reach.
\end{remark}

\section*{What these three have in common}

Ergodic separation, staleness, and intervention are three distinct
ways a system's continuation geometry can be misjudged, by structure,
by elapsed time, and by force, and this chapter has kept them formally
apart on purpose. Collapsing them into one diagnosis, an observer
failing to track reality, would have been easier to write and less
useful to read: each has its own signature, its own proof, and its
own remedy, and only one of the three is fixed by looking again.

\chapter{Investment-Movable Admissibility}

Not every boundary asks the same thing of you. A species cannot
breathe an atmosphere its lungs did not evolve for, and no amount of
committed effort changes that fact; the boundary is simply where it
is, permanently, for reasons that have nothing to do with anyone's
budget. A planet's atmosphere, on the other hand, can be thickened,
warmed, seeded, and made breathable by exactly the kind of committed
effort the first boundary is indifferent to. Both are limits on what
is admissible. Only one of them recedes when you push on it.

Every constraint set $K$ used so far in this book and its predecessor
has been fixed once and for all, a single unmoving boundary between
what a system may become and what it may not. This chapter asks what
changes when that boundary is allowed to move, not arbitrarily, but in
response to something accumulated over time, and does so using
exactly the operator vocabulary Chapter 4 already built, applied to a
kind of object that chapter did not yet ask it to touch.

\section*{Constraint sets are shapes too}

Chapter 4 defined a continuation operator as a Lipschitz map $T :
\mathcal{K}(X) \to \mathcal{K}(X)$ and used it exclusively to
transform reachable sets. Nothing in that definition mentions
reachability. A constraint set $K$, provided it is compact, is simply
another element of $\mathcal{K}(X)$, and there is nothing preventing
an operator from acting on it directly.

\begin{definition}[Investment family]
An investment family is a collection of continuation operators
$\{E_\delta\}_{\delta \geq 0}$ on $\mathcal{K}(X)$ satisfying:
\begin{enumerate}[label=(\roman*)]
\item $E_0 = \mathrm{id}$;
\item $E_{\delta_1 + \delta_2} = E_{\delta_2} \circ E_{\delta_1}$ for
all $\delta_1, \delta_2 \geq 0$ (a one-parameter semigroup);
\item $K \subseteq E_\delta(K)$ for every $K \in \mathcal{K}(X)$ and
$\delta \geq 0$ (expansive).
\end{enumerate}
Given a base constraint set $K_0$, the investment schedule is $K_i :=
E_i(K_0)$ for accumulated investment $i \geq 0$.
\end{definition}

\begin{proposition}[The schedule is monotone]
\label{prop:monotone}
$K_i \subseteq K_j$ whenever $i \leq j$.
\end{proposition}
\begin{proof}
Write $j = i + \delta$ for $\delta = j - i \geq 0$. By the semigroup
property, $K_j = E_j(K_0) = E_\delta(E_i(K_0)) = E_\delta(K_i)$, and by
expansiveness, $K_i \subseteq E_\delta(K_i) = K_j$.
\end{proof}

\begin{remark}
This is the formal content of ``investment never shrinks admissibility'':
once the family satisfies the three axioms above, monotonicity is not
an additional assumption, it is forced. What the axioms do not force,
and what distinguishes one investment family from another, is how
much room a given increment of investment actually buys, which is
exactly the content of the operators $E_\delta$ themselves rather than
of the schedule they generate.
\end{remark}

\section*{Hard and movable boundaries}

\begin{definition}[Hard-excluded set]
$K_\infty := \overline{\bigcup_{i \geq 0} K_i}$, the closure of the
investment orbit. A state $x \in X$ is hard-excluded if $x \notin
K_\infty$; otherwise $x$ is investment-reachable.
\end{definition}

\begin{remark}
The closure is necessary, not decorative: a monotone union of compact
sets need not itself be compact, or even closed, since the union can
approach a boundary point without ever reaching it. $K_\infty$ as
defined here is the smallest closed set containing every stage of the
schedule, which coincides with the plain union exactly when that
union happens to be closed already, and not otherwise.
\end{remark}

\begin{remark}
This makes the opening distinction exact. A species trait sits outside
$K_\infty$ for every conceivable investment family a design could
choose, not merely the one at hand; a planet's habitability threshold
sits inside $K_i$ for $i$ large enough, under an investment family
that models what terraforming can actually do. Whether a given
boundary is hard or movable is not a property of the boundary alone;
it is a property of the boundary together with whichever investment
family is available to act on it, and the same physical limit can be
hard under one family and movable under a richer one.
\end{remark}

\begin{proposition}[Viability is upward-closed in investment]
\label{prop:viab-upward}
If $x \in \operatorname{Viab}_F(K_i)$ for some $i$, then $x \in
\operatorname{Viab}_F(K_j)$ for every $j \geq i$.
\end{proposition}
\begin{proof}
$x \in \operatorname{Viab}_F(K_i)$ gives a trajectory $\gamma \in
\Gamma_{K_i}(x)$ remaining in $K_i$ for all time. By
Proposition~\ref{prop:monotone}, $K_i \subseteq K_j$, so $\gamma$
remains in $K_j$ for all time as well, giving $\gamma \in
\Gamma_{K_j}(x)$ and hence $x \in \operatorname{Viab}_F(K_j)$.
\end{proof}

\begin{remark}
Once a state becomes viable at some investment level, no further
investment can undo that; the set of investment levels at which a
given state is viable, if nonempty, is upward-closed. This is a weak
fact stated for a reason: it is exactly strong enough to justify
treating investment as a one-way ratchet on viability and no stronger,
since Proposition~\ref{prop:viab-upward} says nothing about whether a
state \emph{becomes} viable at any finite level in the first place,
only that viability, once reached, is never lost to further
investment.
\end{remark}

\section*{A worked schedule}

\begin{example}[Habitability threshold]
Let $X = [0,1]$ represent a planet's habitability score, and let
$\tau(i) := \max(\tau_{\min},\, \tau_0 - c i)$ for constants $\tau_0
\in (0,1)$, $c > 0$, and a floor $\tau_{\min} \in [0, \tau_0)$
representing an absolute physical limit, pressure or temperature
bounds no terraforming budget can cross. Define
$K_i := [\tau(i), 1]$.
\end{example}

\begin{proposition}
This schedule arises from a valid investment family, and
$K_\infty = [\tau_{\min}, 1]$, already closed.
\end{proposition}
\begin{proof}
Define $E_\delta$ on all of $\mathcal{K}(X)$ by the dilation
$E_\delta(K) := \{x \in [0,1] : \operatorname{dist}(x,K) \leq c\delta\}
\cap [\tau_{\min}, 1]$, capping at the hard floor. On $[0,1] \subset
\mathbb{R}$ with the standard metric, dilation composes exactly,
$E_{\delta_2}(E_{\delta_1}(K)) = E_{\delta_1+\delta_2}(K)$, since a
point within $c\delta_2$ of a point within $c\delta_1$ of $K$ is
within $c(\delta_1+\delta_2)$ of $K$ by the triangle inequality, and
conversely, giving the semigroup property; $K \subseteq E_\delta(K)$
is immediate, giving expansiveness; and $E_\delta$ is Lipschitz in
Hausdorff distance with constant $1$, a standard property of metric
dilations. Restricted to the orbit generated from $K_0 = [\tau_0,1]$,
$E_\delta([\tau(i),1]) = [\max(\tau_{\min}, \tau(i) - c\delta), 1] =
[\tau(i+\delta), 1]$, matching the schedule's definition exactly. As
$i \to \infty$, $\tau(i)$ reaches $\tau_{\min}$ \emph{exactly} once $i
\geq i_0 := (\tau_0 - \tau_{\min})/c$, not merely in the limit, so
$K_i = [\tau_{\min},1]$ for all $i \geq i_0$, and
$\bigcup_i [\tau(i),1] = [\tau_{\min},1]$ is already closed: for this
particular schedule, the closure in the definition of $K_\infty$ adds
nothing, though it would be needed for a schedule approaching its
floor only asymptotically.
\end{proof}

\begin{remark}
Everything below $\tau_{\min}$ is hard-excluded regardless of budget;
everything in $[\tau_{\min}, \tau_0)$ is investment-reachable but not
part of the original admissible set, exactly the region terraforming
is supposed to buy; and $\tau_{\min}$ itself is where the two kinds of
boundary this chapter opened with meet, the point past which no
further increment of $E_\delta$ does anything at all, since $\tau$ is
already at its floor.
\end{remark}

\section*{An open question}

\begin{openproblem}
Chapter 1 showed that a family of \emph{contracting} operators drives
any starting shape to a unique attractor from outside, in the sense
that repeated application converges in Hausdorff distance regardless
of where it starts. An investment family is expansive rather than
contracting, but it is monotone and, in examples like the one above,
bounded above by a fixed $K^{\mathrm{hard}}$. Whether a bounded,
monotone, expansive operator family converges to a well-defined limit
shape in an analogous sense, a kind of attractor approached from
below rather than surrounded from outside, is not addressed here.
The habitability example converges to $[\tau_{\min},1]$ by direct
computation, but nothing in this chapter establishes that this always
happens, only that it happened once.
\end{openproblem}

\section*{Closing}

Nothing about the operators in this chapter required a new
definition; Chapter 4's continuation operators, applied to constraint
sets instead of reachable sets, were already enough. What was new was
the question, not the machinery, and that is the right way around for
a book that has tried, throughout, to reuse what it already has before
reaching for something else.

\chapter{Semantic Manifolds}

Call a light red and you have already thrown away almost everything
true about it: which fraction of a second into the cycle it is, how
many pedestrians are waiting, whether a request was just registered
or has been pending for a while. None of that survives the word. And
yet the word is not a failure of description, it is the entire point
of having one. A category is supposed to discard detail. The question
this chapter actually has to answer is not whether ``red'' throws
things away, every category does, but whether there is a principled
mathematical object underneath the word at all, or whether ``meaning''
was always going to be a modeling choice dressed up as a discovery.

An earlier sketch of this project wanted to say that a semantic
category is an equivalence class of histories under
$\Gamma(h_1) \cong \Gamma(h_2)$, and stopped there, with $\cong$ never
specified. That is not a theorem with a gap in it. It is a sentence
that has not yet decided what it is about. This chapter decides,
and part of deciding honestly is admitting that the most obvious
candidate is wrong before proposing a better one.

\section*{The obvious candidate is too strict}

The most literal reading of $\Gamma(h_1) \cong \Gamma(h_2)$ available
in this book is equality, $\Gamma(h_1) = \Gamma(h_2)$, which is
exactly $\sim_\Gamma$ from Chapter 2. Nothing new would be needed to
state a theorem this way. Something would be lost by doing so.

\begin{proposition}
Ordinary semantic categories are, in general, coarser than
$\sim_\Gamma$: there exist $h_1, h_2$ commonly called by the same
name with $\Gamma(h_1) \neq \Gamma(h_2)$.
\end{proposition}
\begin{proof}
Chapter 6, Theorem~1: $(0,0)$ and $(0,1)$, both ordinarily described
as ``waiting at the red light,'' satisfy $R_1(0,0) = \{(1,0),(1,1)\}
\neq \{(1,1)\} = R_1(0,1)$.
\end{proof}

\begin{remark}
This was already proved for a different purpose, and it settles the
question immediately: if $\cong$ is read as literal equality, ``red
light'' is not a semantic category in this book's sense at all, since
two states any ordinary speaker would call by that name already fail
to be $\sim_\Gamma$-equivalent. Meaning, whatever it turns out to be,
cannot be this strict without contradicting how the word is actually
used.
\end{remark}

\section*{A structural alternative}

A weaker and more promising idea is that two situations share a
meaning not because their futures are identical, but because their
futures have the same \emph{shape}, related by some symmetry of the
system rather than by outright coincidence.

\begin{definition}[Structural equivalence]
Fix the metric $d((k,r),(k',r')) := d_{\mathrm{cyc}}(k,k') + |r-r'|$
on $X = \{0,1,2,3\} \times \{0,1\}$, where $d_{\mathrm{cyc}}(k,k') :=
\min(|k-k'|, 4-|k-k'|)$ is the cyclic distance on phase. Let $G$ be a
group of isometries of $(X,d)$ acting \emph{equivariantly}: $g(R_T(x))
= R_T(g(x))$ for every $g \in G$ and every $x$. States $x,y$ are
$G$-equivalent, $x \sim_G y$, if $y = gx$ for some $g \in G$.
\end{definition}

\begin{remark}
This is a genuine equivalence relation for the reason a single fixed
map is not: reflexivity, symmetry, and transitivity follow
automatically from $G$ being a group (the identity, inverses, and
closure under composition, respectively), whereas $x \sim_\varphi y
\iff y = \varphi(x)$ for one fixed $\varphi$ need not be symmetric or
transitive unless $\varphi$ itself happens to generate such a group.
The equivariance condition is doing real work too: it is what makes
$\sim_G$ interact sensibly with $\Gamma$ at all, since without it $G$
could relate states with unrelated futures for no reason connected to
the dynamics.
\end{remark}

\begin{example}[An isometry that is not a symmetry]
Let $\varphi_1 : (k,r) \mapsto ((k+1) \bmod 4,\, r)$, the phase shift.
$\varphi_1$ is an isometry of $(X,d)$, since cyclic distance is
translation-invariant. Direct computation gives
$\varphi_1(R_1(k,r)) = R_1(\varphi_1(k,r))$ for $k=0$ and $k=1$, and
$\varphi_1(R_1(k,r)) \neq R_1(\varphi_1(k,r))$ for $k=2$ and $k=3$: at
$k=0$, $\varphi_1(R_1(0,0)) = \{(2,0),(2,1)\} = R_1(1,0)$, matching,
while at $k=2$, $\varphi_1(R_1(2,0)) = \{(0,0),(0,1)\} \neq
\{(0,0)\} = R_1(3,0)$, failing.
\end{example}

\begin{remark}
This computation shows $\varphi_1$ fails equivariance, and the
honest conclusion is narrower than a first pass at this example
suggested. Since equivariance must hold everywhere for $\langle
\varphi_1 \rangle$ to generate a valid $G$, and it demonstrably does
not (it fails at $k=2,3$), $\varphi_1$ cannot be used to build a
nontrivial structural equivalence on this system at all under
Definition 3: the only group of isometries equivariant with this
particular $\Phi$ across the whole space is the trivial one. What the
computation does show, and what is worth keeping, is that $\varphi_1$
is a genuine isometry of the state space that happens to agree with
the dynamics on part of its domain and disagree on the rest, phases
$0,1$ against phases $2,3$, precisely the phases adjacent to
clearing. That is a real fact about where the system's transition
rule does and does not resemble a uniform relabeling, but it is a
fact about a \emph{local}, not global, correspondence, best formalized
by a groupoid of partial symmetries rather than a group if pursued
further, an extension not carried out here. A word like ``waiting''
that treats all four phases alike is, on this reading, not grounded
in even a partial group symmetry of the dynamics; it is grounded in
nothing more than the coarse observation $\Pi$ discussed already, and
the appeal of the phase-shift computation is diagnostic rather than
foundational: it locates precisely where a uniform label stops being
defensible, without supplying the label a genuine structural warrant
anywhere.
\end{remark}

\section*{What a semantic category actually is}

Neither $\sim_\Gamma$ nor $G$-equivalence is what ordinary
language means by a category, and this chapter's honest conclusion is
that nothing in this book's vocabulary should be expected to be. A
semantic category, ``red light,'' ``waiting,'' ``sleeping,'' is
generally an observation map $\Pi$ in the sense of Chapter 6, chosen
for reasons external to continuation fidelity entirely: convenience,
communicability, whatever distinctions happened to be salient to
whoever built the vocabulary. Chapters 2, 3, and 6 did not need this
chapter's help to already supply the complete toolkit for auditing
such a choice once made: whether $\Pi$ is state-sufficient
(Chapter 6), and if not, exactly how much it loses and at which fibers
(Chapter 3). What this chapter adds is a second, independent question
worth asking of the same $\Pi$: does it happen to align with a genuine
equivariant symmetry of the system, in the sense of $G$-equivalence,
or does it merely resemble one locally, as ``waiting'' does on the
traffic light, agreeing with a candidate isometry on two of four
phases while failing the equivariance a genuine group would need
everywhere. A category can be lossy and still correspond to a real,
if partial, structural pattern; it can also be lossy and correspond
to no pattern in the dynamics at all. This book's machinery can tell
these apart even when neither rises to the strength of a genuine
group symmetry. It cannot tell you which category to use, because
that was never a mathematical question in the first place.

\begin{remark}[On the original theorem]
The equivalence class construction this chapter was meant to
formalize, meaning as continuation-equivalence rather than
state-equivalence, survives, but not as a single theorem with one
canonical $\cong$. It survives as a choice between at least two
precise, inequivalent relations, $\sim_\Gamma$ (too strict for
ordinary use) and $G$-equivalence under an equivariant symmetry group
(precise, but demanding enough that this book's own worked example
satisfies it only trivially), together with the recognition that most
categories anyone actually uses are neither, and are auditable against
both without being reducible to either.
\end{remark}

\chapter{Triple Stores and Reachability Closure}

A relational table starts from entities and attaches properties to
them. A graph database of the kind an RDF triple store implements
starts from relations and lets entities emerge as whatever sits at
the ends of them. A single fact, Paris is the capital of France,
becomes an edge, not a row, and a database built entirely out of
edges like this one turns out to behave less like a filing cabinet
and more like something this book has spent its length studying: a
small set of local rules whose repeated application reaches much
further than what was explicitly written down. The temptation is to
call this generation in exactly the sense of Chapter 1's fern. That
temptation should be resisted, and resisting it properly is most of
what this chapter does.

\section*{Reachability, not generation}

\begin{definition}
For entities $X$ and triples $E \subseteq X \times P \times X$ (with
$P$ a set of predicates), the one-step reachability operator is
$\Phi(x) := \{y : (x,p,y) \in E \text{ for some } p \in P\}$, and
$R_T(x) := \bigcup_{n=0}^{T} \Phi^n(x)$ is the set of entities
reachable within $T$ edges.
\end{definition}

\begin{remark}
This is not new. It is the Kauffman correspondence of Chapter 2 of
\emph{Continuation Geometry} \cite{kauffman2000}, already proved there
by induction to satisfy $S_{t+1} = S_t \cup \mathrm{Adj}(S_t)$, applied
here with $\mathrm{Adj}(S_t) := \bigcup_{x \in S_t} \Phi(x)$. That
chapter was careful to distinguish this construction from generation
in the sense of a Hutchinson attractor, and the distinction matters
enough here to restate: $\Phi$ is not, in general, a contraction of
anything. There is no metric under which traversing an edge shrinks
distances, no completeness argument, no Banach fixed point, and
correspondingly no unique limit shape that repeated traversal
converges to. $R_T(x)$ simply grows, a monotone union exactly as
Proposition 1 of \emph{Continuation Geometry}'s Chapter 2 already
proved for this recursion in general. A triple store is an instance
of accumulation, not of the contracting machinery Chapter 1 of this
book built. Calling it generative in that stronger sense would be the
same conflation this book's own first chapter took care to rule out.
\end{remark}

\begin{remark}
What survives of the analogy is real, and worth keeping separate from
what does not. A small, explicit generator set producing a much
larger implicit structure is common to both an IFS attractor and a
triple store's transitive closure; only the mechanism differs, metric
contraction to a fixed point in one case, combinatorial branching
outward from a seed in the other. The next section makes the
branching case precise on its own terms rather than borrowing Chapter
1's vocabulary for it.
\end{remark}

\begin{center}
\begin{tabular}{ll}
\hline
\textbf{This book} & \textbf{RDF / graph databases} \\
\hline
one-step reachability $\Phi$ & edge traversal \\
$R_T(x)$ & $T$-hop neighborhood \\
continuation volume $\operatorname{Vol}_T$ & neighborhood size \\
$\sim_\Gamma$ & indistinguishable neighborhoods \\
continuation quotient & graph compression \\
sufficient statistic (Chapter 5 of \emph{Continuation Geometry}) & context coordinate \\
continuation distance $d_\Gamma$ & embedding similarity \\
\hline
\end{tabular}
\end{center}

\begin{remark}
This table is a reading aid, not a further claim; every row is
established, corrected, or left explicitly open in the sections that
follow, and none of them should be taken as settled by appearing in a
table.
\end{remark}

\section*{How much is implied by how little}

\begin{proposition}[Reachable volume under bounded branching]
\label{prop:volbound}
If $|\Phi(x)| \leq b$ for every $x$, then
$|R_T(x)| \leq \sum_{k=0}^{T} b^k$.
\end{proposition}
\begin{proof}
By induction on $T$. $|R_0(x)| = 1 \leq b^0$. Assume
$|R_T(x)| \leq \sum_{k=0}^T b^k$. Every element of $R_{T+1}(x)$ is
either already in $R_T(x)$ or reached by one further edge from some
element of $\Phi^T(x)$, so
$|R_{T+1}(x)| \leq |R_T(x)| + |\Phi^T(x)| \cdot b \leq |R_T(x)| +
b^{T+1}$ (bounding $|\Phi^T(x)|$ by $b^T$, itself an immediate
consequence of the same bound applied one level down), giving
$|R_{T+1}(x)| \leq \sum_{k=0}^{T+1} b^k$.
\end{proof}

\begin{corollary}
For $n$ entities each with out-degree at most $b$, edge storage is
$O(nb)$, while a single entity's $T$-step reachable volume can be as
large as $\Theta(b^T)$, exceeding total edge storage once
$T > \log_b(n)$ roughly.
\end{corollary}

\begin{remark}
This is the honest version of the compression claim: not that the
store is somehow secretly larger than it appears, but that the
quantity worth measuring, $T$-step reachable volume in the sense of
Chapter 6's $\operatorname{Vol}_T$, specialized here to counting
measure on a discrete graph, grows combinatorially in $T$ while
storage grows linearly in $n$. A social graph storing a hundred
million edges can have individual $T$-step neighborhoods, for modest
$T$, larger than the edge set itself. Nothing about this requires the
store to compute or cache those neighborhoods in advance; $\Phi$
alone, applied repeatedly at query time, is enough.
\end{remark}

\section*{Context coordinates as retained history}

A triple store proper cannot represent two facts about the same
subject and predicate at once. A quad store, storing $(s,p,o,g)$ for
a context or graph label $g$, can. Whether dropping $g$ is a
representational loss or an operational one, in the sense of Chapter
10, depends on which kind of store is being asked.

\begin{example}[Representational: a query drops context]
A quad store holds both $(\text{Alice}, \text{worksAt}, \text{Acme},
2024)$ and $(\text{Alice}, \text{worksAt}, \text{BetaCorp}, 2026)$. A
query that ignores $g$ sees both employers as simultaneously true,
losing the temporal distinction, but the underlying store still holds
both facts and a differently written query recovers them. This is
Chapter 3's fiber loss, not Chapter 11's operational loss: nothing was
destroyed, only misread by one particular query.
\end{example}

\begin{example}[Operational: the schema has no $g$ slot]
A plain triple store with a functional predicate, at most one object
per subject-predicate pair, cannot hold both facts at all; inserting
the second triple overwrites or conflicts with the first. Here the
distinction is gone in the sense Chapter 11 gives that word: no query
against this store, however cleverly written, recovers the discarded
employer, because the schema itself never retained the information
needed to answer it.
\end{example}

\begin{corollary}[Context coordinates are sufficient-statistic augmentation]
\label{cor:context-augmentation}
Let $X$ be the space of $(s,p,o)$ triples and let $g$ range over
contexts. If $\Gamma(s,p,o,g) \neq \Gamma(s,p,o,g')$ for some $s,p,o$
and $g \neq g'$, then the plain-triple observation $\Pi(s,p,o,g) :=
(s,p,o)$ is state-insufficient in the sense of Chapter 6, and the
quad-store state $(s,p,o,g)$ is exactly the augmentation
$x \mapsto (x, \varphi(H_t))$ of Chapter 5 of \emph{Continuation
Geometry}, with $g$ playing the role of the retained coordinate
$\varphi(H_t)$.
\end{corollary}
\begin{proof}
This is Chapter 5 of \emph{Continuation Geometry} instantiated
directly: that construction augments a state $x$ with a summary
$\varphi(H_t)$ of history exactly when $\Gamma$ depends on more than
$x$ alone, and requires nothing beyond checking that dependence, which
is the stated hypothesis $\Gamma(s,p,o,g) \neq \Gamma(s,p,o,g')$
here. No new argument is needed; the triple store's context label is
an instance of an already-proved construction, not a new one.
\end{proof}

\begin{remark}
The graph label $g$, read this way, is a sufficient statistic in the
sense of Chapter 5 of \emph{Continuation Geometry}: it is exactly the
augmented coordinate that construction called for, $x$ silently
becoming $(x, g)$, when a single observable state is not enough to
determine what comes next. A triple store without context coordinates
is a system that has already discarded that augmentation; whether the
discarding is fixable is not a fact about the data remaining somewhere
unread, it is a fact about whether the schema was built to keep it.
\end{remark}

\section*{Meaning, corrected}

\begin{remark}
It is tempting, watching graph embeddings work well in practice, to
write $\mathrm{Meaning}(x) \approx \Gamma(x)$ and stop there. Chapter
9 already showed why the naive version of that claim, meaning as
literal reachable-set equality, is too strict: two entities with
almost, but not exactly, identical neighborhoods would count as
having nothing in common, which is not how embeddings or ordinary
semantic judgment actually behave. What graph embeddings compute is
closer to a distance than an equality, neighborhoods aggregated and
compared by similarity rather than checked for exact coincidence,
which is Chapter 6's $d_\Gamma$, not $\sim_\Gamma$. Two entities being
semantically close because their $T$-step neighborhoods are near in
Hausdorff distance is a defensible, gradeable claim in a way that
Chapter 9 already proved bare equality is not. The correction Chapter
9 made for a four-state traffic light applies without modification to
a hundred-million-edge graph.
\end{remark}

\begin{remark}[The shape of the connection to graph neural networks, and its limit]
The structural form $\mathrm{Embedding}(x) := F(R_T(x))$, for some
aggregation $F$, names the right correspondence: an embedding is a
compressed coordinate on a continuation space, exactly matching this
book's $d_\Gamma$ once $F$ is chosen so that similar aggregated
outputs correspond to nearby reachable sets. What this remark does
not claim is that graph neural networks compute this $F$ mechanically
as stated. Real architectures aggregate learned feature vectors
attached to nodes through differentiable, trained functions, layer by
layer, not a fixed $F$ applied once to the raw reachable set $R_T(x)$;
the correspondence here is a claim about shape, an aggregation of
neighborhood structure into a compressed coordinate, not a claim that
this book's formalism reproduces the mechanics of message passing or
attention. Treating $F(R_T(x))$ as a structural analogy for what an
embedding is doing, rather than a specification of how one is
computed, is the honest version of this connection.
\end{remark}

\section*{What compresses, made exact}

\begin{corollary}[Graph quotient]
\label{cor:graphquotient}
Fix horizon $T$. The relation $x \sim_T y \iff R_T(x) = R_T(y)$ is an
equivalence relation on graph nodes, and the quotient graph it induces
is the coarsest graph preserving every $T$-bounded reachability query
against the original.
\end{corollary}
\begin{proof}
$\sim_T$ is $\sim_\Gamma$ specialized to $\Gamma := R_T$, so it is an
equivalence relation by the definition of Chapter 2. That the induced
quotient is coarsest among those preserving $T$-bounded queries is
Theorem 1 of Chapter 3 (the coarsest zero-loss quotient), instantiated
with $H$ the node set and $\Gamma := R_T$: any coarser identification
merges some pair with $R_T(x) \neq R_T(y)$, which some $T$-bounded
query distinguishes, by definition of $R_T$ as the record of every
node reachable within $T$ edges.
\end{proof}

\begin{remark}
Graph compression by neighborhood equivalence is not analogous to a
continuation quotient. It is one, with no translation required beyond
choosing $\Gamma := R_T$.
\end{remark}

\begin{remark}[An unplanned validation of Chapter 6]
Graph practitioners already draw, informally and by necessity, the
same distinction Chapter 6 built formally: a node's identifier and a
node's neighborhood are different things, and two nodes can share one
while differing in the other. A recommendation system that
deduplicates users by neighborhood similarity, ignoring stated
identity, and a knowledge graph that keeps two identically-labeled
entities apart because their surrounding facts differ, are both
already committed to the claim Chapter 6 had to argue for from first
principles: that an observable label need not be sufficient to
determine what a node's structure actually is. This does not make
Chapter 6's argument redundant. It makes it independently confirmed,
by a field that arrived at the same distinction for entirely practical
reasons and never needed to call it state insufficiency to rely on it.
\end{remark}

\section*{Closing}

A triple store is a real system built, apparently independently of
this book's concerns, around the same primitive this book has spent
its length formalizing: not the state of an entity, but what remains
reachable from it. Where the analogy to Chapter 1's attractors was
tempting and wrong, this chapter said so and pointed to the
correspondence that is actually apt, Chapter 2's accumulation. Where
the analogy to Chapter 9's meaning was tempting and too strong, this
chapter softened it to the distance Chapter 6 already built. What
needed no correction at all, the compression argument and the
context-coordinate reading, is offered as this chapter's genuine
contribution: not a new theorem, but two places where an existing one
turned out to already be about something built for entirely different
reasons a decade before this book was written.

\chapter{Representational Loss versus Operational Loss}

Tell someone that a system ``cleared its memory'' and they will
picture one of two entirely different events, and ordinary language
gives them no way to know which. In the first, a record-keeper simply
starts describing two situations the same way going forward, a
paperwork decision that changes nothing about the world the paperwork
describes. In the second, something is physically erased, and the
capability that depended on it is gone regardless of what anyone
writes down afterward. A summary of this book's own traffic-light
example blurred exactly this line, narrating a change of label as
though it were a change of wiring, and the blur is worth resolving
properly, since the two events have the same size in bits and
opposite consequences in the world.

\section*{Two kinds of forced merge}

\begin{definition}[Representational loss]
For a quotient $q : H \to O$, the representational loss at a fiber $o$
is $\mathcal{L}(o) := \log |C_\Gamma(o)|$, exactly as in Chapter 3:
the log-count of distinct continuation spaces a description choice
groups under one label, with the underlying dynamics $\Phi$ left
untouched.
\end{definition}

\begin{definition}[Operational loss]
For an intervention $T$ (Chapter 7) applied on a set $S$, the
operational loss on $S$ is $\mathcal{L}_{\mathrm{op}}(T,S) := \log
|\{\Gamma(x) : x \in S\}|$, computed on the \emph{pre-intervention}
values of $\Gamma$: the log-count of distinct futures $S$'s members
actually had before $T$ forced them to a common target.
\end{definition}

\begin{remark}
The two definitions look identical, and the resemblance is the point,
not an accident to apologize for. What differs is not the formula but
what it is being asked about. Representational loss is a property of
a description laid over an unchanged system: the pedestrian's request
is still there, still shaping the real controller's real next state,
and only the label an observer uses to talk about it has become too
coarse to reveal that. Operational loss is a property of the system
after something has actually acted on it: the request itself is gone,
and no amount of relabeling recovers it, because there is nothing
left to relabel accurately.
\end{remark}

\section*{What Chapters 2 and 3 actually did}

\begin{proposition}
\label{prop:purelyrep}
Every computation in Chapters 2 and 3 involving the traffic light's
alternative clearing rule used a single fixed $\Phi$ throughout;
$\mathcal{L}_{\mathrm{op}} = 0$ everywhere in that material, since no
intervention was ever applied.
\end{proposition}
\begin{proof}
The alternative rule considered in Chapter 2 (clearing whenever
currently at phase $0$, rather than on arrival) was introduced and
evaluated entirely as a candidate quotient $q$, compared against the
one fixed transition function $\Phi$ given in that chapter's
definition. Chapter 3's loss computation, likewise, computed
$C_\Gamma(o_{\mathrm{bad}})$ using $R_1$ values generated by that same
single $\Phi$. At no point was a second transition function
constructed, nor was any operator applied to a state to force its
subsequent behavior; every quantity computed was a property of a
label held up against one unchanging system.
\end{proof}

\begin{remark}
This is worth stating as a proposition rather than a footnote,
because the natural way to narrate the same material out loud tends
to drift toward the operational reading without anyone intending it
to. Saying a design ``forcibly changes state $A$ into state $B$'' is
a perfectly natural way to describe a coarse label, but it is also,
read literally, a description of an intervention, and the two readings
support different conclusions. Under the representational reading,
nothing about the intersection's actual behavior was ever at stake;
the pedestrian's request was tracked correctly by $\Phi$ the entire
time, and the only failure was in how a hypothetical observer using
the bad label would have talked about it. Proposition~\ref{prop:purelyrep}
is the precise statement that this is, in fact, all that was shown.
\end{remark}

\section*{The operational twin}

\begin{example}[The same bit, actually erased]
Let $T'$ be the operator sending $(0,1) \mapsto (0,0)$ (and fixing
every other state), a genuine intervention in the sense of Chapter 7:
it forces the physical request bit to $0$ at the moment of
intervention, not merely relabels the situation.
\end{example}

\begin{proposition}
$\mathcal{L}_{\mathrm{op}}(T', \{(0,0),(0,1)\}) = \log 2$, exactly the
value Chapter 3 computed for the representational case.
\end{proposition}
\begin{proof}
The pre-intervention values are $\Gamma(0,0) = R_1(0,0) =
\{(1,0),(1,1)\}$ and $\Gamma(0,1) = R_1(0,1) = \{(1,1)\}$, already
established as distinct in Chapter 6. So $|\{\Gamma(0,0),\Gamma(0,1)\}|
= 2$, giving $\mathcal{L}_{\mathrm{op}} = \log 2$.
\end{proof}

\begin{remark}
The number is identical to Chapter 3's one bit, and this identity of
magnitude is exactly what makes the two cases easy to conflate and
important to keep separate. Under $T'$, a pedestrian who ran up and
pressed the button while the light sat at phase $0$ genuinely loses
the system's acknowledgment of that press: the controller's actual
future behavior is now the same whether or not the button was pressed,
because the operator overwrote the state that would have carried the
difference forward. This is not a labeling failure a better catalog
could fix. The information a better label would have preserved is no
longer there to preserve.
\end{remark}

\section*{Why the difference is not academic}

\begin{remark}
A purely representational loss is repaired for the cost of a better
label, since Proposition~\ref{prop:purelyrep}'s kind of failure never
touched the underlying system; recomputing $\sim_\Gamma$ correctly and
using it instead costs nothing beyond noticing the old label was
wrong. An operational loss, once incurred, is repaired only by
whatever it takes to restore the erased distinction physically, if
that is even possible, since Chapter 7's machinery already
established that an intervention forcing $\Gamma(T(x))$ to a fixed
target is not, in general, invertible by any further relabeling.
Mistaking the second kind of failure for the first is not a rounding
error. It is the difference between believing a problem is free to
fix and not noticing that it no longer can be.
\end{remark}

\chapter{Closing}

The word in this book's title was chosen before most of the book was
written, which is the ordinary way titles get chosen and also, on
reflection, a reasonable place to end by checking. A fern grown from
four contractions is generative in the most literal sense available:
nothing was accumulated, everything was produced by rules acting on
their own output. Not everything that came after Chapter 1 earned the
word as cleanly, and an honest closing chapter has to say so rather
than let the title do the arguing for chapters that moved on to
something else.

\section*{What was established}

Stated plainly: this book proved the Continuation Attractor Theorem
in full, a genuine Hutchinson-and-Banach result applied to the
Hausdorff-metric space of reachable sets already built in the first
book, with a worked, verified Cantor-set instance. It defined the
continuation quotient without circularity, having first caught and
discarded an earlier attempt that assumed its own conclusion, and
proved its universal property and its invariance property from that
definition alone. It proved that the coarsest quotient losing no
continuation-relevant information is $\sim_\Gamma$ itself, tying
directly back to the sufficient-statistic material of the first
book's Chapter 5 rather than introducing a competing notion. It
computed, exactly and by verified arithmetic rather than by hand-wave,
where a real toy system's design was provably correct and where a
superficially similar alternative would have provably destroyed
information, and it caught a genuine numerical error in its own first
attempt at that computation before publishing the corrected version.
It placed generation and collapse in one vocabulary,
a continuation operator, and proved that vocabulary closed under
composition regardless of where an operator sits on the genuinely
separate axes of contraction and injectivity, catching and correcting
an overclaim of its own that had compressed those axes into one
number. It proved a real Collage Theorem, three lines
long and built entirely from material the first chapter had already
established, and used it to prove that every compact target is
approximable, while keeping that result honestly separated from the
much harder conjecture it is often mistaken for. It resolved an
explicitly flagged gap, state insufficiency, without inventing a new
toy system to do it, by observing an already-built one differently.
It gave staleness and invariant-block separation different proofs and
different remedies rather than one gesture toward both, reusing a
Gr\"onwall estimate from the first book's own appendix rather than
deriving a new one. It showed that investment-movable admissibility
needed no new operator definition at all, only a decision to point
Chapter 4's machinery at constraint sets instead of shapes. It
resolved its own explicitly flagged undefined symbol, settling on two
precise and mutually irreducible notions of continuation-equivalence
rather than one overloaded one, and said plainly when the answer
reduced to material already on hand. It found, in an industrial graph
database built for entirely unrelated reasons, a genuine instance of
this book's own primitive, and was honest that the instance fit one
chapter's vocabulary well and another's not at all, correcting a
tempting but wrong analogy to generation before it could take root.
And, close to the end, it caught
a conflation in a plausible-sounding retelling of its own second
chapter and gave the conflation a name and a formal treatment rather
than a footnote.

\section*{What was not}

The Universal Continuation Generation conjecture, named precisely in
Chapter 5, was not proved, and this book's own methods do not reach
it; the gap between existence of an approximating family and
existence of a single universal one was made exact, not closed.
Whether a genuinely dimensional, rather than merely cardinality-based,
obstruction to finite sufficient statistics exists remains exactly
where the first book left it, untouched here. The question of whether
a bounded, monotone, expansive operator family converges to a
well-defined limit from below, Chapter 8's open problem, was
demonstrated once and proved never. Whether any operator actually
named Collapse, Refuse, or Project elsewhere in this corpus is
literally, technically a contraction or a quotient in the senses
built here was repeatedly flagged as an empirical question about
those specific constructions and was never once investigated. The
nested-constraint account of agency raised during this book's
planning was set aside for lack of any formal target this vocabulary
could give it, and it should stay set aside until one is found rather
than be forced into whichever chapter seemed most convenient. And no
chapter here built a general repair theory for operational loss
analogous to the first book's account of repair as capture; Chapter
10 named the distinction between a loss that can be undone by
relabeling and a loss that cannot, but said nothing about how, or
whether, the second kind is ever actually recovered.

\section*{Was generative the right word}

Read the table of contents cold and the honest pattern is this: two
chapters, the first and part of the eighth, are generative in the
sense the title promises, operators building structure up from
nothing but their own repeated application. A comparable number,
the second, third, sixth, seventh, ninth, tenth, and eleventh, are
about the opposite motion: quotients collapsing distinctions, losses
computed and located, insufficiencies made explicit, failure modes
kept formally apart, conflations caught and corrected, an industrial
example's own generative-sounding analogy corrected on arrival. The
fourth chapter
is the hinge, though its own closing claim had to be corrected before
it could bear the weight this paragraph puts on it: generation and
collapse are not two regions of one Lipschitz parameter, a single
number never separated them, only injectivity did that, independently
of contraction. What the fourth chapter actually earned is narrower
and still real: both motions share one vocabulary, continuation
operators, and are checkable by the same closure property regardless
of which region of that richer, multi-axis space they occupy. The
book was never really about generation any more than it was about
collapse. It was about the vocabulary both turned out to need.

\section*{A fourth thread the count above simplifies}

The two-bucket count above, generative against everything else, is
honest as far as it goes, but it is not the finest division available,
and pretending otherwise would repeat exactly the kind of premature
compression this book has spent its length arguing against. Chapters
7 and 8 do not answer the question the rest of the book keeps
returning to: whether two histories, states, or descriptions can be
treated as the same without losing anything continuation-relevant.
Invariant decomposition is not a quotient in that sense at all; it
says nothing about identifying $x$ with $y$, only that some region of
the state space is causally sealed off from the rest, a fact about a
set, not a relation on a pair. Staleness and intervention are not
equivalence questions either; they concern what happens to a fixed
piece of knowledge, or a fixed state, as time elapses or as force is
applied, a question about change over time rather than about which
things coincide. Chapter 8's investment operators inherit the same
character: a boundary moving as effort accumulates is a claim about a
schedule, not about which states are interchangeable.

What Chapters 7 and 8 share, and what generation, quotienting, and
sufficiency alike do not quite have room for, is a concern with time
and force themselves: how a fact about a system degrades, is
confined, or is overridden as the clock runs or as effort is spent,
independent of whether two things were ever the same to begin with.
This book is better described by four threads than by two:
generation, quotienting, sufficiency, and the dynamics of
admissibility and knowledge under time and intervention. Naming the
fourth one now, rather than folding it into whichever of the first
three it resembled most, is the same discipline this book applied to
its own Lipschitz overclaim in Chapter 4 and its own undefined
equivalence in Chapter 9: a count that comes out even is not evidence
the count was right.

\begin{remark}
If this book has a real center of gravity, it is not the fern and it
is not the confluence of rivers. It is the discovery, made honestly
and only in retrospect, that a project which set out to explain how
futures get built spent most of its actual length explaining how they
get lost, mislabeled, confined, staled, and mistaken for one another,
and that neither half makes sense as a theory on its own. A theory of
generation with nothing to say about loss would not have been able to
tell a correct design from an incorrect one in Chapter 2. A theory of
loss with nothing to say about generation would have had no attractor
to measure the loss against in the first place. The title stays,
because the object it names, a continuation operator, really does
generate both halves from one definition. But anyone reading this book
for a theory of how futures are built should know, going in, that they
are about to spend at least as much time learning how futures are
taken away.
\end{remark}

\chapter*{References}
\addcontentsline{toc}{chapter}{References}

\begin{thebibliography}{99}
\bibitem{banach1922}S.\ Banach, ``Sur les op\'erations dans les
ensembles abstraits et leur application aux \'equations
int\'egrales,'' Fundamenta Mathematicae, 3 (1922).
\bibitem{falconer2003}K.\ Falconer, \emph{Fractal Geometry:
Mathematical Foundations and Applications}, 2nd ed., Wiley, 2003.
\bibitem{barnsley1988}M.\ Barnsley, \emph{Fractals Everywhere},
Academic Press, 1988.
\bibitem{hutchinson1981}J.\ E.\ Hutchinson, ``Fractals and
self-similarity,'' Indiana University Mathematics Journal, 30 (1981).
\bibitem{stone1948}M.\ H.\ Stone, ``The generalized Weierstrass
approximation theorem,'' Mathematics Magazine, 21 (1948).
\bibitem{cybenko1989}G.\ Cybenko, ``Approximation by superpositions
of a sigmoidal function,'' Mathematics of Control, Signals, and
Systems, 2 (1989).
\bibitem{aubinframkowska1990}J.-P.\ Aubin, H.\ Frankowska,
\emph{Set-Valued Analysis}, Birkh\"auser, 1990.
\bibitem{kauffman2000}S.\ Kauffman, \emph{Investigations}, Oxford
University Press, 2000.
\end{thebibliography}

\end{document}
